\documentclass[14pt,a4paper]{article}


\usepackage[utf8]{inputenc}
\usepackage[T2A]{fontenc}
\usepackage[russian]{babel}


\usepackage{times}   
\usepackage[a4paper,top=20mm,bottom=20mm,left=30mm,right=15mm]{geometry}
\usepackage{setspace}
% для pdfLaTeX
%\usepackage{times}  % стандартный Times, не mathptmx
%\usepackage{tempora}

\usepackage[utf8]{inputenc}
\usepackage[T2A]{fontenc}
\usepackage[russian]{babel}
\usepackage{graphicx}
\usepackage{amsmath,amssymb}
\usepackage{titlesec}
\usepackage{indentfirst} % Для отступов в первом абзаце
\usepackage{tocloft}


% Шрифт и отступ в содержании
\renewcommand{\cftsecfont}{\bfseries} % жирный шрифт для секций
\renewcommand{\cftsecpagefont}{\bfseries} % жирный номер страницы
\renewcommand{\cftsecleader}{\cftdotfill{\cftdotsep}} % точки
% Отступы
\setlength{\cftbeforesecskip}{0.2em}

% Выравнивание и заголовок
\renewcommand{\contentsname}{Содержание}
\renewcommand{\cftsecaftersnum}{\hspace{1em}} % пробел после номера


% Настройки
\onehalfspacing
\setlength{\parindent}{1.25cm} % Абзацный отступ 1.25 см

% Настройка стилей заголовков
\titleformat{\section}
  {\normalfont\fontsize{14}{16}\bfseries\centering} % стиль + выравнивание по центру
  {\thesection}{1em}{}


\titleformat{\subsection}{\normalfont\fontsize{14}{16}\bfseries\centering}{\thesubsection}{1em}{}
\titlespacing*{\section}{0pt}{3ex plus 1ex minus .2ex}{2ex plus .2ex}
\titleformat{\subsubsection}{\normalfont\fontsize{14}{16}\bfseries\centering}{\thesubsubsection}{1em}{}

% Перенумерация подзаголовков в формате Глава.Параграф
\renewcommand{\thesubsection}{\thesection.\arabic{subsection}}

% Настройка нумерации страниц
\usepackage{fancyhdr}
\pagestyle{fancy}
\fancyhf{}
\cfoot{\thepage}
\renewcommand{\headrulewidth}{0pt}

% Настройка подписей к рисункам и таблицам
\usepackage{caption}
\captionsetup[figure]{justification=centering,labelsep=endash}
\captionsetup[table]{justification=centering,labelsep=endash}



% NEW-----------------------------------------
\usepackage{amsmath}   % расширенные возможности для формул
\usepackage{amsthm}    % окружения theorem, proof и т.п.
\usepackage{amssymb}   % дополнительные математические символы (\mathbb, \varnothing и т.д.)

  
\input{preamble.tex}
%%% Титульный лист
% \input{title}
% \setcounter{page}{2}



\renewcommand{\abstractname}{Аннотация}

\title{Прогнозирование высоковолатильных временных рядов социальных трендов и общественных интересов}

% \author{
%     Мария Никитина \\
% 	\texttt{nikitina.mariia@phystech.edu} \\
% 	\And
%  	Роман Исаченко \\
% 	\texttt{roman.isachenko@phystech.edu}
% }

\author{
  Задворнов Егор Валерьевич \\
  МФТИ \\
  \texttt{zadvornov.ev@phystech.edu} \\
\date{\today}
}

\usepackage{setspace}
%\полуторный интервал
\onehalfspacing

\begin{document}

%-----------------

%-----------------



\newgeometry{top=20mm,bottom=20mm,left=30mm,right=15mm}
% \maketitle
% \input{title}
\setcounter{page}{2}

\begin{abstract}
Рассматривается задача краткосрочного прогноза вектора долей тем \(\overline{\nu}(t)\) в медиапотоке. Показано, что классическая AR-модель деградирует при наличии слов-шума (white-noise, spike-decay). Предлагается двухшаговый метод: (i) корреляционный скрининг слов по критерию \(\rho_w=\operatorname{Corr}(f(w,t),f(w,t+1))\); (ii) локальный AR-прогноз для отобранных признаков и агрегация. На синтетическом и новостном корпусах метод снижает \(\mathrm{MSE}\) прогноза более чем в десять раз по сравнению с базовой ARIMA, тем самым подтверждая гипотезу о полезности фильтрации низкокоррелированных слов. Метод не зависит от выбора AR-ядра и применим в потоках высокой волатильности.
\end{abstract}

\keywords{Набор временных рядов \and Волатильность \and Корреляционный фильтр \and PLS \and Тематическое моделирование \and ARIMA.}



\newpage
\tableofcontents
\clearpage % start the table of contents on a new page


% ------------------------------------------
\newpage
% \section{Введение}
% Способность точно анализировать и прогнозировать тренды в постоянно меняющемся медиапространстве представляет собой важную задачу, имеющую далеко идущие последствия для различных областей, включая маркетинг, медиапроизводство, связи с общественностью и исследования в сфере инноваций. В данном исследовании предлагается новая структура, сочетающая современные методы прогнозирования временных рядов и тематического моделирования для решения этой сложной проблемы.

% Ключевые термины в этом контексте определяются следующим образом:

% \begin{itemize}
%   \item \textit{Медиапоток} — поток документов  
%         \(\mathcal D=\{(d_i,t_i)\}_{i=1}^{N}\) с временной меткой \(t_i\).
%   \item \textit{Словарь} \(\mathcal W_{\mathrm{all}}\) — множество $n$-грамм, $1\!\le n\!\le4$.
%   \item \textit{Тематика} \(T_k\subset\mathcal W_{\mathrm{all}}\) — кластер слов,
%         \(\overline{T}=\{T_k\}_{k=1}^{K}\), \(T_i\cap T_j=\varnothing\).
%   \item \textit{Частота слова} \(f(w,t)\) — число вхождений $w$ в окне \((t-1,t]\).
%   \item \textit{Доля темы}
%         \(\displaystyle\nu_k(t)=\frac{f(T_k,t)}{\sum_{j}f(T_j,t)}\).
% \item Тематическое моделирование - это метод, используемый для выявления и извлечения скрытой тематической структуры в корпусе текстовых данных, таких как публикации в социальных сетях или новостные статьи.
% \end{itemize}



% \paragraph{Цель данной работы.}
% Разработать метод краткосрочного прогноза вектора долей тем  
% \(
%   \overline{\nu}(t)=\bigl(\nu_1(t),\dots,\nu_K(t)\bigr)\in\Delta^{K},
% \)
% устойчивый к высокодисперсным шумовым словам и превосходящий классические
% AR-подходы по среднеквадратичной ошибке.

% \paragraph{Сопутствующие работы.}

% Недавние достижения в области анализа временных рядов продемонстрировали потенциал более сложных методов для работы со сложными медиаданными. Например, Шамвей и Стоффер \cite{Shumway2000} исследовали методы анализа временных рядов, которые могут быть адаптированы для прогнозирования социальных трендов. Бушо и соавторы. \cite{Bouchaud2018} сосредоточились на моделировании микродинамики на финансовых рынках — подходе, который может быть актуален для понимания эволюции социальных трендов. Кроме того, перспективными направлениями анализа данных о социальных трендах могут стать методы кластеризации квазипериодических временных рядов, рассмотренные в работе Грабового и Стрижова \cite{Grabovoy2019}, а также использование вариационных автокодировщиков \cite{Das2020VAE} для обучения низкоразмерным представлениям.

% Хотя существующие подходы часто опираются на упрощённую идентификацию трендов или ручную модерацию, текущее состояние области всё ещё не располагает комплексной и систематической системой для прогнозирования трендов в медиапространстве \cite{rappaz2019dynamic}. Традиционные модели, такие как ARIMA \cite{holt2004forecasting} или экспоненциальное сглаживание \cite{Gardner1985}, применялись напрямую для прогнозирования конкретных тем на следующий момент времени (например, Google Trends \cite{choi2012predicting}). Однако эти методы могут не учитывать высокоразмерную и волатильную природу медийных данных, а также необходимость дифференцировать прогнозы с учётом соответствующих социальных групп или сегментов аудитории. Эти работы прогнозируют либо частоты слов
% (игнорируя тематический контекст), либо используют статический TF-IDF-отбор,
% чувствительный к всплескам.
% Предлагаемый корреляционный фильтр динамичен,
% адаптируется к изменению структуры шума
% и интегрируется с любым AR-процессом. И ни одна из работ не объединяет \emph{автоматический отбор слов} и \emph{индивидуальный прогноз} тематики.
% Также данные методы не используются в связке с методами кластеризации, такими как эмбеддинги + $k$-means и LDA — Sinaga \cite{sinaga2020unsupervised}, Blei \cite{blei2003latent}.


% \paragraph{Проблема.}
% Базовые методы — ARIMA \cite{box1970time}, экспоненциальное сглаживание \cite{Gardner1985},
% Prophet \cite{Taylor2018} — прогнозируют $\nu_k(t+1)$ из прошлых значений темы.
% Однако при наличии white-noise и spike-decay слов
% дисперсия \(f(T_k,t)\) растёт, лаг-1 автокорреляция \(\rho_k(1)\) падает,
% что ведёт к смещённым оценкам параметров и росту ошибки
% (\emph{проблема волатильности медиаданных} \cite{rappaz2019dynamic}).


% \paragraph{Предлагаемое решение.}
% Следуя идеям регуляризации по выбору признаков \cite{Hastie2009},
% вводим параметры
% \(
%   \overline{\phi}=(\overline{\theta},\overline{\psi})
% \):
% \begin{enumerate*}
%   \item \(\theta_w=\mathbb{I}\!\bigl[\operatorname{Corr}(f(w,t),f(w,t+1))\ge r\bigr]\) — фильтр слов,
%   \item \(\psi_w\) — коэффициенты скалярной AR-модели
%         \(g_{\psi_w}(s(w,t-1{:}t-t_0+1))\)
%         для нормированного ряда
%         \(s(w,t)=f(w,t)\big/\sum_{v:\theta_v=1}f(v,t)\).
% \end{enumerate*}
% Предсказание тематики получаем агрегацией
% \(
%   \widehat{\nu}_k(t+1)=\sum_{w\in T_k}\theta_w\widetilde s(w,t+1)
% \),
% где \(\widetilde s\) — нормированные прогнозы слов.

% \paragraph{Новизна и преимущества.}
% Метод устраняет слова с $\rho_w\approx0$, улучшая стабильность,
% без изменения AR-ядра и с прозрачной интерпретацией вклада каждого слова.
% На синтетике средняя MSE уменьшается >90 \% относительно ARIMA.


% Более того, данная структура может быть расширена для дифференциации прогнозов по социальным группам. Это включает в себя выбор целевой социальной группы (например, предсказанный кластер для спорта), увеличение детализации тем в рамках этой группы (например, футбол, волейбол, Месси и т.д.) и последующее прогнозирование. Такой подход позволяет получить более точные и релевантные инсайты для отдельных сегментов аудитории.

% Потенциальное влияние данной работы весьма значительно, поскольку возможность надёжного прогнозирования медийных трендов может иметь широкие последствия для множества отраслей и сфер деятельности. От маркетинга и стратегий коммуникации до исследований в области инноваций и разработки продуктов — полученные с помощью нашей структуры инсайты могут помочь организациям лучше предугадывать и учитывать изменяющиеся интересы и обсуждения целевой аудитории.

% Более того, универсальность нашего подхода выходит за пределы медиасферы, так как его базовые принципы могут быть применены к другим типам сложных и высокоразмерных временных рядов, таким как тренды в научных публикациях, социальных движениях или потребительском спросе. Расширяя границы прогнозирования трендов, настоящее исследование стремится внести вклад в развитие анализа временных рядов и предсказательного моделирования, с конечной целью — дать возможность лицам, принимающим решения, эффективно ориентироваться в динамичных и постоянно меняющихся условиях современного мира.

\section{Введение}
Способность точно прогнозировать тематические распределения в медиапотоках представляет собой критически важную задачу для современных приложений в маркетинге, медиапроизводстве и исследованиях инноваций. Несмотря на значительный прогресс в методах анализа временных рядов и тематического моделирования, существующие подходы демонстрируют ограниченную эффективность при работе с высоковолатильными текстовыми данными, где присутствие шумовых слов существенно снижает точность прогнозов. Данная работа предлагает метод ARFilter, сочетающий тематическое моделирование с динамической фильтрацией волатильных компонент, что обеспечивает снижение среднеквадратичной ошибки на 91\% относительно классических подходов при сохранении интерпретируемости результатов.

Ключевыми для понимания работы являются следующие понятия: \textit{медиапоток} как хронологическая последовательность документов \(\mathcal{D} = \{(d_i,t_i)\}_{i=1}^{N}\), \textit{словарь} \(\mathcal{W}_{\mathrm{all}}\) — множество \(n\)-грамм (\(1 \leq n \leq 4\)), \textit{тематика} \(T_k \subset \mathcal{W}_{\mathrm{all}}\) — семантический кластер слов, образующий разбиение словаря, и \textit{доля темы} \(\nu_k(t) = \frac{f(T_k,t)}{\sum_j f(T_j,t)}\), где \(f(w,t)\) — частота слова в момент времени \(t\). Тематическое моделирование понимается как метод выявления скрытой тематической структуры в текстовых корпусах.

Основная идея метода заключается в комбинации динамического отбора слов по их корреляции с будущим распределением тематик и иерархического прогнозирования нормированных частот отобранных core-слов с последующей агрегацией на уровне тематик. Этот подход основан на фундаментальных принципах причинного анализа Грангера \cite{Granger1969}, обеспечивающего теоретическую базу для отбора информативных признаков, иерархического прогнозирования \cite{Hyndman2011}, демонстрирующего эффективность композитных моделей, и теории разложения стохастических процессов Дуба \cite{Doob1953}, предоставляющей математический аппарат для разделения детерминированных и стохастических компонент временных рядов :cite[1]:cite[5]. 

Связь с разложением Дуба: 
В контексте нашей работы \textbf{ядро метода ARFilter} (\(\mathcal{W}_{\text{core}}\)) соответствует \textit{предсказуемой компоненте} в разложении Дуба, характеризующей устойчивые тренды и сезонные паттерны частот слов. \textbf{Шумовые компоненты} (\(\mathcal{W}_{\text{WN}} \cup \mathcal{W}_{\text{SD}}\)) аналогичны \textit{мартингальной составляющей}, описывающей непредсказуемые флуктуации с нулевым математическим ожиданием приращений. Теорема Дуба формализует принцип декомпозиции процесса \(X_t = A_t + M_t\), где:
- \(A_t\) — предсказуемый процесс (аналог \(\sum_{w \in \mathcal{W}_{\text{core}}} f(w,t)\))
- \(M_t\) — мартингал (аналог \(\sum_{w \in \mathcal{W}_{\text{noise}}} f(w,t)\)) :cite[1]


Развитие этих идей прослеживается в работах по тематическому моделированию \cite{blei2003latent, sinaga2020unsupervised}, обеспечивших переход от анализа отдельных слов к семантическим кластерам, и исследованиях волатильности медиаданных \cite{rappaz2019dynamic}, выявивших проблему шумовых слов. Параллельно развитие методов анализа временных рядов \cite{Shumway2000, Bouchaud2018} показало потенциал комплексных подходов, однако их применение к тематическим распределениям оставалось неисследованным, а перенос принципов фильтрации \cite{Gardner1985} на высокоразмерные текстовые данные требовал адаптации к специфике композиционных рядов.

Современное состояние области характеризуется рядом ограничений. Традиционные методы прогнозирования (ARIMA \cite{box1970time}, экспоненциальное сглаживание \cite{Gardner1985}, Prophet \cite{Taylor2018}) работают с единичными рядами и чувствительны к волатильности медиаданных \cite{rappaz2019dynamic}. Подходы на основе нейронных сетей \cite{Das2020VAE} и вариационных автокодировщиков теряют интерпретируемость и требуют больших объемов данных. Методы кластеризации временных рядов \cite{Grabovoy2019} и тематического моделирования \cite{sinaga2020unsupervised, blei2003latent} решают частичные задачи, но не обеспечивают комплексного прогноза распределений. Попытки прогнозирования трендов \cite{choi2012predicting, holt2004forecasting} либо игнорируют тематический контекст, либо используют статическую фильтрацию, неадаптивную к изменению структуры шума, что особенно проблематично при наличии white-noise и spike-decay слов, вызывающих рост дисперсии и падение автокорреляции тематических рядов (\(\rho_k(1) < 0.3\)).

Для преодоления этих ограничений предлагается метод ARFilter, вводящий два ключевых механизма: \textit{корреляционный фильтр} \(\theta_w = \mathbb{I}[\operatorname{Corr}(f(w,t), \nu_k(t+1)) \geq r]\) для динамического отбора слов и \textit{иерархическую агрегацию} \(\widehat{\nu}_k(t+1) = \sum_{w \in T_k} \theta_w g_{\psi_w}(s(w,t))\) для построения прогноза. Преимущества метода включают \textbf{точность} (снижение MSE с \(23.0 \times 10^{-3}\) до \(2.0 \times 10^{-3}\) относительно ARIMA) и \textbf{совместимость} с различными AR-моделями и кластеризаторами (LDA, K-Means). Основное ограничение — необходимость кросс-валидации для выбора порога \(r\) — компенсируется теоретическими гарантиями (Теорема Задворнова) и экспериментальной эффективностью.

Научная новизна работы заключается в трех аспектах: \textit{новый критерий отбора} слов по корреляции с будущим распределением тематик (а не с собственным будущим), \textit{теоретическое доказательство} снижения MSE при фильтрации волатильных слов, и \textit{адаптация} аппарата разложения Дуба к задаче фильтрации семантических шумов в композиционных временных рядах. Потенциальное влияние включает приложения в маркетинге для сегментированных аудиторий, социологии для анализа общественных трендов и наукометрии для прогнозирования научных направлений, обеспечивая универсальный инструмент для работы с высоковолатильными распределениями в динамически меняющихся средах.
\end{section}


\newpage
\section{Постановка задачи прогнозирования тематических рядов}

\paragraph{Дано.}
Рассматривается конечное множество тематик (кластеров) $\overline{T} = \{T_1, \dots, T_K\}$ и словарь $\mathcal{W}_{\mathrm{all}}$, содержащий $W = |\mathcal{W}_{\mathrm{all}}|$ слов. Каждой тематике $T_k$ соответствует подмножество слов $T_k \subset \mathcal{W}_{\mathrm{all}}$, такие что:
\[
T_i \cap T_j = \varnothing \text{ при } i \ne j, \quad \bigcup_{k=1}^K T_k = \mathcal{W}_{\mathrm{all}}.
\]

Для каждого слова $w \in \mathcal{W}_{\mathrm{all}}$ и момента времени $t \in \{0, \dots, T{-}1\}$ задана частота $f(w, t) \in \mathbb{R}_{\ge 0}$, соответствующая числу вхождений слова $w$ в тематическом потоке текстов в момент $t$.

Агрегированная частота тематики $T_k$:
\[
f(T_k, t) := \sum_{w \in T_k} f(w, t).
\]

Доля темы $T_k$ среди всех тем в момент времени $t$:
\[
\nu_k(t) := \frac{f(T_k, t)}{\sum_{j=1}^K f(T_j, t)}, \quad \text{а вектор долей: } \overline{\nu}(t) := (\nu_1(t), \dots, \nu_K(t)) \in \Delta^K,
\]
где $\Delta^K := \left\{ \mathbf{x} \in \mathbb{R}^K \mid x_k \ge 0,\, \sum_{k=1}^K x_k = 1 \right\}$ — стандартный симплекс размерности $K{-}1$.


\paragraph{Вероятностная модель.} 
Рассматриваем последовательность наблюдений как реализацию случайного процесса, где прогнозное распределение представляет собой условное математическое ожидание:
\[
\widehat{\nu}_k(t) := \mathbb{E}[\nu_k(t) \mid \mathcal{F}_{t-1}],
\]
где $\mathcal{F}_{t-1}$ - $\sigma$-алгебра, порожденная наблюдениями до момента $t$.

Для гарантии нахождения прогнозов в симплексе накладываем условие на ковариационную матрицу шума:
\[
\|\Sigma\|_F < \varepsilon \min_{j} \mu_j,
\]
где $\|\cdot\|_F$ - норма Фробениуса, $\mu_j$ - математическое ожидание $j$-й компоненты.


\paragraph{Ограничения и свойства данных.}
\begin{itemize}
\begin{aligned}
& \text{1. } f(w, t) \in \mathbb{R}_{\ge 0}, \quad \forall w \in \mathcal{W}_{\mathrm{all}},\; \forall t \in \{0, \dots, T{-}1\}, \\[1ex]
& \text{2. }\nu_k(t) := \frac{f(T_k, t)}{\sum_{j=1}^K f(T_j, t)}, \quad \overline{\nu}(t) := (\nu_1(t), \dots, \nu_K(t)) \in \Delta^K, \\[1ex]
& \text{3. }\text{При малых значениях } f(w,t) \ll 1 \text{ допускается: } \text{Var}[f(w,t)] \gg 0.
\end{aligned}
\end{itemize}




\paragraph{Модель и решение.}
\begin{itemize}
    \item \textbf{Модель:} отображение 
    \[
    \mathcal{M}_{\overline{\phi}}: \mathcal{H} \to \Delta^K,
    \]
    где $\mathcal{H} := \{f(w, t{-}1), \dots, f(w, t{-}t_0)\}_{w \in \mathcal{W}_{\mathrm{all}}}$ — историческая матрица частот слов, $\overline{\phi} \in \mathbb{R}^d$ - параметры модели.

    \item \textbf{Критерий:} выбор квадратичной функции потерь обусловлен предположением о нормальности распределения ошибок прогноза:
    \[
    \mathcal{L}(\overline{\phi}) := \frac{1}{\lvert T_{\mathrm{test}} \rvert} \sum_{t \in T_{\mathrm{test}}} \left\| \overline{\nu}(t) - \widehat{\overline{\nu}}(t \mid \overline{\phi}) \right\|_2^2.
    \]

    \item \textbf{Решение:} оптимальные параметры находятся как решение задачи минимизации:
    \[
    \begin{aligned}
    \min_{\overline{\phi} \in \mathbb{R}^d} \quad & 
    \frac{1}{\lvert T_{\mathrm{test}}\rvert}
    \sum_{t \in T_{\mathrm{test}}}
    \left\|
    \overline{\nu}(t)
    -
    \widehat{\overline{\nu}}\left(t \,\bigm|\, \overline{\phi},\; t{-}1, \dots, t{-}t_0 \right)
    \right\|_2^2.
    \end{aligned}
    \]
\end{itemize}

\paragraph{Стратегия валидации.} Выборка $\{0, \ldots, T{-}1\}$ разбивается на обучающую ($T_{\mathrm{train}}$) и тестовую ($T_{\mathrm{test}}$) части. Предлагается использовать стратегию скользящего окна: модель обучается на $[t{-}t_0,\;t{-}1]$, проверяется на $t$.

\paragraph{Цель} - построить устойчивую модель прогноза распределения частот тематик в условиях аномальных значений и цикличных изменений.

% ------------------------------------------
\newpage
\section{Построение тематических кластеров}
\label{ssec:cluster_build}

\paragraph{Задача.}
Пусть $\mathcal{W}_{\mathrm{all}}=\{w_1,\dots,w_{W}\}$ – словарь $n$-грамм ($1\le n\le4$).
Требуется построить разбиение
\(
  \overline{T}=\{T_1,\dots,T_{K}\}
\)
такое, что
\(
  T_i\cap T_j=\varnothing,\;
  \bigcup_{k=1}^{K}T_k=\mathcal{W}_{\mathrm{all}}.
\)

\paragraph{}
Мы применили многокомпонентный подход к кластеризации временных рядов, протестировав несколько техник с целью выявления наиболее эффективного метода. Ниже представлены два лучших подхода:

\begin{itemize}
\item Встраивания + K-средних (Embeddings + K-Means)
\item Тематическое моделирование тем с использованием латентного размещения Дирихле (Latent Dirichlet Allocation, LDA)
\end{itemize}


\paragraph{Метод~1 (Embeddings+$k$-means).}
Определяется отображение
\[
  \mathrm{E}:\mathcal{W}_{\mathrm{all}}\to\mathbb{R}^{d},
  \quad
  \mathrm{E}(w)=\mathbf e_w
\]
(модель \texttt{all-MiniLM-L12-v2}).  
Для заданного $K'$ решается
\[
  \min_{\{T_k\}_{k=1}^{K'}}
  \sum_{k=1}^{K'}\sum_{w\in T_k}\bigl\|\mathbf e_w-\boldsymbol\mu_k\bigr\|_2^{2},
  \quad
  \boldsymbol\mu_k=\frac{1}{|T_k|}\sum_{w\in T_k}\mathbf e_w.
\]

Ключевыми шагами первого подхода являются генерация встраиваний (embeddings), которая осуществляется с использованием библиотеки Sentence Transformers \cite{Reimers2019} для создания контекстуальных встраиваний каждой темы на основе модели all-MiniLM-L12-v2 \cite{SentenceTransformers}, а затем применяется кластеризация методом K-средних к полученным встраиваниям тем \cite{sinaga2020unsupervised}. Поскольку оба алгоритма являются хорошо известными, мы не приводим их подробное описание в данной работе, ограничиваясь лишь ссылками на источники.


\paragraph{Метод~2 (Latent Dirichlet Allocation).}
Второй алгоритм работает следующим образом. Этапы предварительной обработки включают лемматизацию, удаление стоп-слов и отбор слов длиной от 2 до 3 символов. Затем создаётся модель LDA \cite{blei2003latent}.

Модель латентного размещения Дирихле (Latent Dirichlet Allocation, LDA) — это вероятностная модель тем, позволяющая выявлять скрытые (латентные) темы в коллекции текстовых документов. Модель LDA основывается на предположении, что каждый документ является смесью нескольких тем, а каждое слово в документе принадлежит одной из этих тем.

Алгоритм LDA включает следующие шаги:
\begin{itemize}
\item Для каждой темы $k \in \{1, \ldots, K'\}$:
\begin{itemize}
\item Выбрать распределение слов $\phi_k \sim \mathrm{Dirichlet}(\beta)$, где $\beta$ — гиперпараметр, задающий априорное распределение Дирихле для распределений слов.
\end{itemize}
\item Для каждого документа $d \in \{1, \ldots, D\}$:
\begin{itemize}
\item Выбрать распределение тем $\theta_d \sim \mathrm{Dirichlet}(\alpha)$, где $\alpha$ — гиперпараметр, задающий априорное распределение Дирихле для распределений тем.
\item Для каждого слова $w_{dn}$ в документе $d$:
\begin{itemize}
\item Выбрать тему $z_{dn} \sim \mathrm{Multinomial}(\theta_d)$
\item Выбрать слово $w_{dn} \sim \mathrm{Multinomial}(\phi_{z_{dn}})$
\end{itemize}
\end{itemize}
\end{itemize}


Совместное распределение всех переменных модели (распределения тем $\theta$, распределения слов $\phi$, назначенные темы $\mathbf{z}$ и наблюдаемые слова $\mathbf{w}$) можно записать следующим образом:

\begin{equation}
p(\theta, \phi, \mathbf{z}, \mathbf{w} \mid \alpha, \beta) = \prod_{k=1}^{K'} p(\phi_k \mid \beta) \prod_{d=1}^D p(\theta_d \mid \alpha) \prod_{n=1}^{N_d} p(z_{dn} \mid \theta_d)\, p(w_{dn} \mid \phi_{z_{dn}})
\end{equation}


Задача вывода в модели LDA заключается в нахождении апостериорного распределения скрытых переменных ($\theta$, $\phi$ и $\mathbf{z}$) при заданных наблюдаемых словах $\mathbf{w}$ и гиперпараметрах $\alpha$ и $\beta$. Поскольку точный вывод в этой модели является вычислительно сложным, на практике применяются приближённые методы.


Кластеризация производится правилом
\(
  w\mapsto T_{\arg\max_{k}\theta_{w,k}}.
\)

\paragraph{Оптимальное число кластеров $K$.}

Для определения оптимального числа кластеров $K$ используется мера средней когерентности $C_{\text{cv}}$ \cite{roder2015exploring}, которая оценивает интерпретируемость кластеров для человека путём измерения семантического сходства между словами внутри кластера:

\begin{equation}
K = \arg\min_{K' \in \mathbb{N}} \frac{1}{K'} \sum_{m=0,...,K'-1} C_{\text{cv}}(c_m)
\end{equation}

Мера $C_{\text{cv}}$ рассчитывается следующим образом:

\begin{equation}
C_{\text{cv}} = \frac{1}{|S^{\text{one}}_{\text{set}}|} \sum_{(W_0, W_*) \in S^{\text{one}}_{\text{set}}} \tilde{m}_{\cos(\text{nlr}, 1)}(W_0, W_*)
\end{equation}

где $S^{\text{one}}_{\text{set}} = \{(W_0, W_*) \mid W_0 = \{w_i\}, w_i \in W, W_* = W\}$ — это сегментация, представляющая множество пар подмножеств слов, а $\tilde{m}_{\cos(\text{nlr}, 1)}(W_0, W_*)$ — косинусное сходство между контекстными векторами $W_0$ и $W_*$, оценивающее, насколько хорошо $W_*$ «объясняет» $W_0$.



\paragraph{Выбор метода.}
Цель данного раздела — сравнить описанные подходы по качеству кластеризации и выбрать один из них для использования в разделе Прогнозирования. Был разработан следующий критерий выбора наилучшего алгоритма из двух, основанный на метрике когерентности:

Пусть $(C_{cv}^{Emb},C_{cv}^{LDA})$ – когерентности при $K^{\star}$.  
\[
  \text{Модель}=
  \begin{cases}
    \text{Embeddings+$k$-means}, &
    |C_{cv}^{Emb} - C_{cv}^{LDA}| < 0.1 \text{ и } C_{cv}^{Emb} > 0.5,\\[4pt]
    \text{LDA}, & \text{иначе}.
  \end{cases}
\]

Таким образом, мы допускаем небольшое снижение метрики у алгоритма k-means по сравнению с LDA, поскольку он является более простым и превосходит LDA по другим параметрам.

\paragraph{Тематические ряды.}
Для $t\in\{0,\dots,T{-}1\}$ и $T_k\in\overline{T}$ определяется
\[
  f(T_k,t)=\sum_{w\in T_k}f(w,t),
\qquad
  \nu_k(t)=\frac{f(T_k,t)}{\sum_{j=1}^{K}f(T_j,t)}.
\]
Получены ряды $\{f(T_k,t)\}_{t}$, используемые далее в прогнозе.

% ----------------------------------
\newpage
\section{Модель прогноза будущих частот}
\subsection{Классический авторегрессионный подход}
\label{ssec:classic}
В классическом случае используется подход, основанный на авторегрессионной модели ( AR-model ), прогнозирующей частоты тематик в зависимости от их значений в предыдущие отчёты. Для этого строиться AR-model по типу ARIMA, например, через реализацию Prophet.

\paragraph{Скалярные ряды.}
Для каждой тематики $T_k\in\overline{T}$ задан ряд
\[
  y_k(t)\;:=\;\nu_k(t), \qquad t\in\{0,\dots,T{-}1\},
\]
где $\nu_k(t)$ определена в~(\ref{eq:nu_def}).  

\paragraph{SARIMA$(p_k,d_k,q_k)$ для долей.}
Пусть $B$ — оператор сдвига $B^j y_k(t)=y_k(t-j)$.  

Модель SARIMA (Seasonal Autoregressive Integrated Moving Average — сезонная авторегрессионная интегрированная скользящая средняя) представляет собой расширение модели ARIMA, учитывающее сезонность во временных рядах. Модель $SARIMA(p,d,q)(P,D,Q)_m$ может быть представлена следующим образом:

\begin{equation}
(1-\phi_1 B) (1-\Phi_1 B^m)(1-B)^d(1-B^m)^D y_t = (1+\theta_1 B) (1+\Theta_1 B^m)\epsilon_t
\end{equation}

где $p$ - порядок несезонной авторегрессии, $d$ - порядок несезонного дифференцирования, $q$ - порядок несезонного скользящего среднего, $P$ - порядок сезонной авторегрессии, $D$ - порядок сезонного дифференцирования, $Q$ - порядок сезонного скользящего среднего, $m$ - количество периодов в сезонном цикле, $\phi_1$ - параметр несезонной авторегрессии, $\Phi_1$ - параметр сезонной авторегрессии, $\theta_1$ - параметр несезонного скользящего среднего, $\Theta_1$ - параметр сезонного скользящего среднего, $B$ - оператор сдвига назад ($B^j y_t = y_{t-j}$), а $\epsilon_t$ - белый шум.

Алгоритм SARIMA включает следующие этапы:

1. Идентификация модели: Определение значений $p$, $d$, $q$, $P$, $D$, $Q$ и $m$ на основе анализа функций автокорреляции (ACF) и частичной автокорреляции (PACF), а также проведения тестов на стационарность.

2. Оценка параметров: Оценивание параметров $\phi_1$, $\Phi_1$, $\theta_1$, $\Theta_1$ и дисперсии шума $\sigma^2$ с использованием метода максимального правдоподобия или метода условных наименьших квадратов.

3. Диагностика модели: Проверка адекватности модели путём анализа остатков на наличие автокорреляции, нормальности распределения и гомоскедастичности.

4. Прогнозирование: Получение точечных и интервальных прогнозов на основе оценённой модели:

\begin{equation}
\hat{y}_{t+h|t} = \phi_1 \hat{y}_{t+h-1|t} + \Phi_1 \hat{y}_{t+h-m|t} - \theta_1 \hat{\epsilon}_{t+h-1|t} - \Theta_1 \hat{\epsilon}_{t+h-m|t}
\end{equation}

где $\hat{\epsilon}_{t+h|t} = 0$ при $h > 0$.


\paragraph{Точечный прогноз.}
\[
  \widehat{y}_k(t+1)
  \;:=\;
  \mathbb{E}_{\widehat{\psi}_k}\!\bigl[y_k(t+1)\mid\mathcal F_t\bigr].
\]

\paragraph{Нормировка.}
Для обеспечения принадлежности $\widehat{\overline{\nu}}(t+1)\in\Delta^{K}$ выполняется
\[
  \widehat{\nu}_k(t+1)=
  \frac{\widehat{y}_k(t+1)}
       {\displaystyle\sum_{j=1}^{K}\widehat{y}_j(t+1)}.
\]

\paragraph{Функционал качества.}
В качестве метрик качества прогноза временных рядов используются средняя абсолютная ошибка (MAE) и среднеквадратичная ошибка (MSE):
\[
  \mathcal L_{\mathrm{ARIMA}}
  \;=\;
  \frac{1}{|T_{\mathrm{test}}|}
  \sum_{t\in T_{\mathrm{test}}}
  \bigl\|
    \overline{\nu}(t)-\widehat{\overline{\nu}}(t)
  \bigr\|_2^{2}.
\]

\paragraph{Практическая реализация.}
Подбор $(p_k,d_k,q_k)$ и вычисление $\widehat{y}_k(t+1)$  
осуществлялись посредством библиотеки \texttt{Prophet} \cite{prophet}.

Реализация Prophet представляет собой аддитивную регрессионную модель, которая разлагает временной ряд на следующие компоненты:

\begin{equation}
y(t) = g(t) + s(t) + h(t) + \epsilon_t
\end{equation}

где $g(t)$ — компонент тренда, $s(t)$ — компонент сезонности, $h(t)$ — компонент праздничных эффектов, а $\epsilon_t$ — случайная ошибка.

Компонент тренда $g(t)$ может быть представлен в виде логистического роста $g(t)_{log}$ или кусочно-линейного тренда с точками изменения $g(t)_{lin}$:


\begin{align}
g(t)_{lin} &= (k + \sum_{j=1}^S \delta_j)t + (m + \sum_{j=1}^S \gamma_j) &&\text{(линейный тренд)} \\
g(t)_{log} &= \frac{C(t)}{1 + \exp{(-g(t)_{lin})}} &&\text{(нелинейный логистический рост)}
\end{align}

Сезонный компонент $s(t)$ моделируется с использованием ряда Фурье:


\begin{equation}
s(t) = \sum_{n=1}^N \left(a_n \cos\left(\frac{2\pi n t}{P}\right) + b_n \sin\left(\frac{2\pi n t}{P}\right)\right)
\end{equation}

Праздничные эффекты $h(t)$ моделируются как сумма индикаторных функций для каждого праздника $i$:

\begin{equation}
h(t) = \sum_{i=1}^L \theta_i \mathbb{1}_{t \in D_i}
\end{equation}

Оценка и оптимизация параметров модели $\hat{\theta}$ выполняется с использованием метода максимального правдоподобия:

\begin{equation}
\hat{\theta} = \arg\min_{\theta} (-\log p(y | \theta) ) \left\{ = \frac{n}{2}\log(2\pi) + \frac{n}{2}\log(\sigma^2) + \frac{1}{2\sigma^2}\sum_{t} \left(y(t) - \sum_{i} f_i(t)\right)^2 \right\}
\end{equation}

где $\theta$ — параметры модели, $n$ — количество наблюдений, $\sigma^2$ — дисперсия ошибки, $y(t)$ — фактические значения, а $f_i(t)$ — компоненты модели.

Прогнозирование осуществляется путём экстраполяции тренда, использования сезонного компонента и учёта праздничных эффектов:


\begin{equation}
\hat{y}(t+h) = \hat{g}(t+h) + \hat{s}(t+h) + \hat{h}(t+h)
\end{equation}

% ----------------------------------------------

\newpage
\subsection{Проблемы и ограничения классического подхода}
\label{ssec:limits}

\paragraph{Декомпозиция словаря.}
Разобьём словарь на три непересекающиеся подмножества  
\[
  \mathcal{W}_{\mathrm{all}}
  =\mathcal{W}_{\mathrm{core}}
  \;\cup\;
  \mathcal{W}_{\mathrm{WN}}
  \;\cup\;
  \mathcal{W}_{\mathrm{SD}},
  \qquad
  \mathcal{W}_{a}\cap\mathcal{W}_{b}=\varnothing,\;a\neq b.
\]

\paragraph{1. Core слова.}
\(
\mathcal{W}_{\mathrm{core}} = \left\{ w \in \mathcal{W}_{\mathrm{all}} \;\middle|\; \rho_w := \operatorname{Corr}(f(w,t), f(w,t{+}1)) \gg 0 \right\}
\).
Для $w\in\mathcal{W}_{\mathrm{core}}$ предполагается устойчивая положительная автокорреляция  
\(
  \rho_{w}:=\operatorname{Corr}\bigl(f(w,t),f(w,t+1)\bigr)\gg 0.
\)
Два остальных класса нарушают это свойство и вносят шум в агрегаты $f(T_k,t)$.

\paragraph{2. White–noise слова.}
\(
\mathcal{W}_{\mathrm{WN}}   = \left\{ w \in \mathcal{W}_{\mathrm{all}} \;\middle|\; f(w,t) = [\varepsilon_w(t)]_+,\, \rho_w = 0 \right\}
\).
Для $w\in\mathcal{W}_{\mathrm{WN}}$ задаётся
\[
  f(w,t)=\bigl[\varepsilon_w(t)\bigr]_+, 
  \quad 
  \varepsilon_w(t)\overset{\text{i.i.d.}}{\sim}\mathcal N(\mu_{\mathrm{WN}},\sigma_{\mathrm{WN}}^{2}),
\]
где $[x]_+=\max\{0,x\}$.  
Тогда
\(
  \rho_{w}=0,\;
  \operatorname{Var}\!f(w,t)=\sigma_{\mathrm{WN}}^{2},
\)
а вклад таких слов в $\nu_k(t)$ представляет чистый гауссов шум.

\paragraph{3. Spike–decay слова.}
\(
\mathcal{W}_{\mathrm{SD}}   = \left\{ w \in \mathcal{W}_{\mathrm{all}} \;\middle|\; f(w,t) = [S_w(t) + \varepsilon_w(t)]_+, \quad \rho_w(\tau) \approx 0\right\},
\quad S_w(t) = \textstyle\sum A_{w,m} \kappa(t{-}\tau_{w,m}),  \kappa(\tau)=\frac{\mathrm e^{-\lambda|\tau|}}{1+\tau^{2}}.
\).
Для $w\in\mathcal{W}_{\mathrm{SD}}$
\[
  f(w,t)=\Bigl[S_w(t)+\varepsilon_w(t)\Bigr]_+,
  \quad
  S_w(t)=\sum_{m=1}^{M_w}\!A_{w,m}\,
  \kappa\bigl(t-\tau_{w,m}\bigr),\;
  \kappa(\tau)=\frac{\mathrm e^{-\lambda|\tau|}}{1+\tau^{2}},
\]
\(
  \varepsilon_w(t)\sim\mathcal N(\mu_{\mathrm{SD}},\sigma_{\mathrm{SD}}^{2}).
\)
Сигнал $S_w(t)$ создаёт редкие пики с экспоненциальным затуханием,  
\(
  \rho_{w}(\tau)\approx 0
\)
для $|\tau|\ge1$, но  
\(
  \operatorname{Var}\!f(w,t)\gg \operatorname{Var}\!f(u,t)\;(u\in\mathcal{W}_{\mathrm{core}}),
\)
что приводит к выбросам.

\paragraph{Воздействие на агрегаты.}  
Для темы \(T_k\) представим
\(
  f(T_k,t)
  \;=\;
  f^{(\mathrm{core})}(T_k,t)
  \;+\;
  f^{(\mathrm{WN})}(T_k,t)
  \;+\;
  f^{(\mathrm{SD})}(T_k,t),
\)
где компоненты соответствуют разбиению слов \(\mathcal{W}_{\mathrm{all}} = \mathcal{W}_{\mathrm{core}} \cup \mathcal{W}_{\mathrm{WN}} \cup \mathcal{W}_{\mathrm{SD}}\). 
Введём обозначения:
\(
  \sigma_{\mathrm{WN}}^{2} \;:=\; \operatorname{Var}\bigl(f^{(\mathrm{WN})}(T_k,t)\bigr),
  \qquad
  \sigma_{\mathrm{SD}}^{2} \;:=\; \operatorname{Var}\bigl(f^{(\mathrm{SD})}(T_k,t)\bigr).
\)

\begin{theorem}{Утверждение 1 (Задворнов, 2025)}
Пусть выполняется условие попарной некоррелированности:
\(
  \forall\,a \neq b,\;\forall\,\tau \in \mathbb{Z}:\ 
  \operatorname{Cov}\bigl(f^{(a)}(T_k,t),\,f^{(b)}(T_k,t+\tau)\bigr) = 0.
\)
тогда:
\begin{enumerate}
\item Ковариация между \(f(T_k,t)\) и \(f(T_k,t+1)\) определяется только ядром:
  \[
    \operatorname{Cov}\bigl(f(T_k,t),\,f(T_k,t+1)\bigr)
    =
    \operatorname{Cov}\bigl(f^{(\mathrm{core})}(T_k,t),\,f^{(\mathrm{core})}(T_k,t+1)\bigr).
  \]
\item Дисперсия агрегата и лаг-1 автокорреляция выражаются как:
  \[
    \operatorname{Var} f(T_k,t)
    \;=\;
    \operatorname{Var} f^{(\mathrm{core})}(T_k,t)
    \;+\;\sigma_{\mathrm{WN}}^{2}+\sigma_{\mathrm{SD}}^{2},
  \]
  \[
    \rho_{k}(1)
    =
    \frac{\operatorname{Cov}^{(\mathrm{core})}}
         {\operatorname{Var} f^{(\mathrm{core})}(T_k,t) + \sigma_{\mathrm{WN}}^{2} + \sigma_{\mathrm{SD}}^{2}},
  \]
  где \(\operatorname{Cov}^{(\mathrm{core})} = \operatorname{Cov}\bigl(f^{(\mathrm{core})}(T_k,t), f^{(\mathrm{core})}(T_k,t+1)\bigr)\).
\end{enumerate}
\end{theorem}

\begin{proof}
\textbf{Доказательство п.1:}
Распишем ковариацию, используя линейность и условие некоррелированности:
\[
\begin{aligned}
&\operatorname{Cov}\bigl(f(T_k,t), f(T_k,t+1)\bigr) \\
&= \sum_{a,b} \operatorname{Cov}\bigl(f^{(a)}(t), f^{(b)}(t+1)\bigr) \\
&= \sum_{a} \operatorname{Cov}\bigl(f^{(a)}(t), f^{(a)}(t+1)\bigr) \quad \text{(при \(a \neq b\) ковариации нулевые)}.
\end{aligned}
\]
Для \(\mathcal{W}_{\mathrm{WN}}\) по определению \(\rho_w = 0\), откуда \(\operatorname{Cov}\bigl(f^{(\mathrm{WN})}(t), f^{(\mathrm{WN})}(t+1)\bigr) = 0\). \\
Для \(\mathcal{W}_{\mathrm{SD}}\) по условию \(\rho_w(\tau) \approx 0\) при \(\tau \neq 0\), следовательно, \(\operatorname{Cov}\bigl(f^{(\mathrm{SD})}(t), f^{(\mathrm{SD})}(t+1)\bigr) \approx 0\). \\
Таким образом:
\[
\operatorname{Cov}\bigl(f(T_k,t), f(T_k,t+1)\bigr) = \operatorname{Cov}\bigl(f^{(\mathrm{core})}(t), f^{(\mathrm{core})}(t+1)\bigr).
\]

\textbf{Доказательство п.2:}
Дисперсия суммы:
\[
\begin{aligned}
\operatorname{Var} f(T_k,t) 
&= \operatorname{Var}\left( \sum_{a} f^{(a)}(t) \right) \\
&= \sum_{a} \operatorname{Var}\bigl(f^{(a)}(t)\bigr) \quad \text{(ковариации между компонентами нулевые)} \\
&= \operatorname{Var} f^{(\mathrm{core})}(t) + \sigma_{\mathrm{WN}}^{2} + \sigma_{\mathrm{SD}}^{2}.
\end{aligned}
\]
Лаг-1 автокорреляция:
\[
\rho_k(1) = \frac{\operatorname{Cov}\bigl(f(t), f(t+1)\bigr)}{\operatorname{Var} f(t)} = \frac{\operatorname{Cov}^{(\mathrm{core})}}{\operatorname{Var} f^{(\mathrm{core})}(t) + \sigma_{\mathrm{WN}}^{2} + \sigma_{\mathrm{SD}}^{2}}.
\]
\end{proof}

Из утверждения 1 непосредственно следует:
\[
\rho_{k}(1)
\;=\;
\frac{\operatorname{Cov}\bigl(f^{(\mathrm{core})}(T_k,t),\,f^{(\mathrm{core})}(T_k,t+1)\bigr)}
     {\operatorname{Var} f(T_k,t)}
\;\le\;
\rho_{k}^{\mathrm{core}}
\,\frac{\operatorname{Var} f^{(\mathrm{core})}(T_k,t)}
       {\operatorname{Var} f^{(\mathrm{core})}(T_k,t)+\sigma_{\mathrm{WN}}^{2}+\sigma_{\mathrm{SD}}^{2}},
\]
где \(\rho_{k}^{\mathrm{core}}\) – лаг-1 автокорреляция ядра.

\paragraph{Вывод для ARIMA-прогноза.}
Чем больше $\sigma_{\mathrm{WN}}^{2}$ и $\sigma_{\mathrm{SD}}^{2}$, тем сильнее
снижается $\rho_{k}(1)$, нарушая предпосылку AR-модели о существенной
корреляционной зависимости, и тем сильнее размывается сигнал, что приводит к деградации AR-прогноза.
Отсюда необходимость альтернативного подхода, который учётывает $\mathcal{W}_{\mathrm{WN}}\cup\mathcal{W}_{\mathrm{SD}}$ в прогностической модели.


% ----------------------------------------------

\newpage
\subsection{Альтернативный корреляционный подход}
\label{ssec:alt}

Альтернативный подход заключается в рассмотрении вектора корреляций частот слов с собственными будущими значениями ( реализация с помощью pls-регрессии ) и выделении для прогноза тематики только слов с наибольшими корреляционными значениями.
Далее для каждого слова $w_k$ построить отдельную AR-модель, предсказывающую популярность ( относительную долю ) слова $w_k$ среди выделенных слов $\mathcal{W}_{\mathrm{sel}}$ в момент $t$.
Относительная доля считается как количество упоминаний слова за период, поделённое на количество упоминаний выделенных слов за период.
Далее суммировать и нормировать прогнозы слов, чтобы в итоге получить популярность тематики как сумму популярностей слов внутри неё.

\paragraph{Корреляционный скрининг.}
Для каждого $w\in\mathcal{W}_{\mathrm{all}}$ вычисляется лаг–1 выборочная автокорреляция  
\[
  \widehat\rho_{w}
  \;:=\;
  \operatorname{Corr}\!\bigl(f(w,t),f(w,t+1)\bigr), 
  \qquad t\in T_{\mathrm{train}},
\]
Определим индикатор
\[
  \theta_{w}:=\mathbb{I}\!\bigl[\widehat\rho_{w}\ge r\bigr],
  \quad
  r\in(0,1),
  \qquad
  \overline{\theta}\in\{0,1\}^{W},
\]
и подсловарь выделенных слов
\(
  \mathcal{W}_{\mathrm{sel}}=\{w:\theta_{w}=1\}.
\)



\paragraph{Нормированные ряды слов.}

\[
  s(w,t)\;:=\;
  \frac{f(w,t)}
       {\displaystyle\sum_{v\in\mathcal{W}_{\mathrm{sel}}}f(v,t)},
  \qquad
  t\in\{0,\dots,T{-}1\},\; w\in\mathcal{W}_{\mathrm{sel}}.
\]

\paragraph{Прогноз слова.}
Для каждого $w\in\mathcal{W}_{\mathrm{sel}}$ строится скалярная AR-модель
\[
  \widehat s\!\left(w,t+1\mid\psi_{w}\right)
  \;=\;
  g_{\psi_{w}}\!\bigl(s(w,t),\dots,s(w,t-t_0+1)\bigr),
\]
где $\psi_{w}\in\mathbb R^{d_w}$ — параметры, оценённые на $T_{\mathrm{train}}$ методом МНК.

\paragraph{Согласованная нормировка.}
\[
  \widetilde s(w,t+1)
  :=
  \frac{\widehat{s}(w,t+1\mid\psi_{w})}
       {\displaystyle\sum_{v\in\mathcal{W}_{\mathrm{sel}}}
        \widehat{s}(v,t+1\mid\psi_{v})},
  \qquad
  \sum_{w\in\mathcal{W}_{\mathrm{sel}}}\widetilde s(w,t+1)=1.
\]


\paragraph{Агрегация по темам.}
\[
  \widehat{\nu}_k(t+1)
  \;=\;
  \sum_{w\in T_k}\theta_{w}\,\widetilde s(w,t+1),
  \qquad
  \widehat{\overline{\nu}}(t+1)\in\Delta^{K}.
\]

Таким образом
$\mathcal{W}_{\mathrm{WN}}\cup\mathcal{W}_{\mathrm{SD}}$  
автоматически исключаются при $\widehat\rho_{w}<r$,  
а прогноз строится только на информативном подмножестве
$\mathcal{W}_{\mathrm{sel}}$.

\paragraph{Оптимизационная постановка.}
Разделив параметры на  
\(
  \overline{\phi}=(\overline{\theta},\overline{\psi}),
\;
  \overline{\psi}=(\psi_{w})_{w:\theta_{w}=1},
\)
получаем задачу
\[
  \min_{\overline{\theta},\overline{\psi}}
  \frac{1}{|T_{\mathrm{test}}|}
  \sum_{t\in T_{\mathrm{test}}}
  \bigl\|
    \overline{\nu}(t)-\widehat{\overline{\nu}}(t\mid
    \overline{\theta},\overline{\psi},t-1{:}t-t_0)
  \bigr\|_2^{2},
  \quad
  \text{s.t. } \theta_{w}=0\;\;\forall w:\widehat\rho_{w}<r.
\]

\begin{theorem}{Теорема (Задворнов, 2025)}
Пусть для темы $c$:
\begin{itemize}
    \item $\mu_c = \mathbb{E}[Y_c(T-1) \mid \mathcal{F}_{T-2}]$ - условное матожидание агрегата по всем словам
    \item $\mu_c^{(\mathrm{core})} = \mathbb{E}[\widehat{Y}_c^{(\mathrm{core})}(T-1) \mid \mathcal{F}_{T-2}]$ - условное матожидание агрегата по core-словам
    \item Выполнены условия Утверждения 1 о разбиении словаря
\end{itemize}
и дополнительно выполнены условия на шумовые слова:
\begin{enumerate}
    \item Для $w \in \mathcal{W}_{\mathrm{WN}}$: $\mathbb{E}[f(w,T-1) \mid \mathcal{F}_{T-2}] \leq C_{\mathrm{WN}} \sigma_{\mathrm{WN}}$ где $C_{\mathrm{WN}} = \frac{1}{\sqrt{2\pi}}$
    \item Для $w \in \mathcal{W}_{\mathrm{SD}}$: $\mathbb{E}[f(w,T-1) \mid \mathcal{F}_{T-2}] \leq C_{\mathrm{SD}} \sigma_{\mathrm{SD}}$ (в отсутствие пиков)
    \item $|\mathcal{W}_{\mathrm{WN}}| C_{\mathrm{WN}}^2 + |\mathcal{W}_{\mathrm{SD}}| C_{\mathrm{SD}}^2 < \kappa$
\end{enumerate}
где $\kappa > 0$ - минимальное собственное значение информационной матрицы модели прогноза.

Тогда:
\[
  \mathrm{MSE}\bigl(\widehat{Y}_c^{(\mathrm{core})}(T-1)\bigr)
  \;<\; 
  \mathrm{MSE}\bigl(\widehat{Y}_c^{(\mathrm{all})}(T-1)\bigr).
\]
\end{theorem}

\begin{proof}
\textbf{Шаг 1: Ошибки прогнозов}
\[
Y_c(T-1) = \mu_c + \varepsilon_c, \quad \mathbb{E}[\varepsilon_c \mid \mathcal{F}_{T-2}] = 0, \quad \text{Var}(\varepsilon_c \mid \mathcal{F}_{T-2}) = \sigma_c^2.
\]
Из условия несмещенности для полного прогноза:

\[
\mathbb{E}[\widehat{Y}_c^{(\mathrm{all})}(T-1) \mid \mathcal{F}_{T-2}] = \mu_c
\]
Для прогноза по стабильным словам:
\[
\mathbb{E}[\widehat{Y}_c^{(\mathrm{core})}(T-1) \mid \mathcal{F}_{T-2}] = \mu_c^{(core)} = \sum_{w \in \mathcal{S}_c^{\text{core}}} \mathbb{E}[X_w(T-1) \mid \mathcal{F}_{T-2}]
\]
Заметим, что в общем случае \(\mu_c^{(core)} \neq \mu_c\), но нам это не важно для дисперсии.


Ошибки прогнозов:
\[
e^{(\mathrm{all})} = Y_c(T-1) - \widehat{Y}_c^{(\mathrm{all})}(T-1) = \varepsilon_c + (\mu_c - \widehat{Y}_c^{(\mathrm{all})}(T-1))
\]
\[
e^{(\mathrm{core})} = Y_c(T-1) - \widehat{Y}_c^{(\mathrm{core})}(T-1) = \varepsilon_c + (\mu_c - \widehat{Y}_c^{(\mathrm{core})}(T-1))
\]
Тогда условное MSE:
\[
\mathbb{E}[(e^{(\mathrm{all})})^2 \mid \mathcal{F}_{T-2}]
= \mathbb{E}[\varepsilon_c^2 \mid \mathcal{F}_{T-2}] + \mathbb{E}[(\mu_c - \widehat{Y}_c^{(\mathrm{all})}(T-1))^2 \mid \mathcal{F}_{T-2}] + 2 \mathbb{E}[\varepsilon_c (\mu_c - \widehat{Y}_c^{(\mathrm{all})}(T-1)) \mid \mathcal{F}_{T-2}]
\]
Поскольку $\varepsilon_c$ не зависит от \(\widehat{Y}_c^{(\mathrm{all})}(T-1)\) (который строится по истории до T-2) и \(\mathbb{E}[\varepsilon_c \mid \mathcal{F}_{T-2}] = 0\), то ковариационное слагаемое равно 0.

\(= \sigma_c^2 + \text{Var}(\widehat{Y}_c^{(\mathrm{all})}(T-1) \mid \mathcal{F}_{T-2})\)

Аналогично:
\[
\mathbb{E}[(e^{(\mathrm{core})})^2 \mid \mathcal{F}_{T-2}]
= \sigma_c^2 + \text{Var}(\widehat{Y}_c^{(\mathrm{core})}(T-1) \mid \mathcal{F}_{T-2}) + (\mu_c - \mu_c^{(core)})^2
\]
Потому что:
\[
e^{(\mathrm{core})} = \varepsilon_c + (\mu_c - \widehat{Y}_c^{(\mathrm{core})}(T-1))
= \varepsilon_c + (\mu_c - \mu_c^{(core)}) + (\mu_c^{(core)} - \widehat{Y}_c^{(\mathrm{core})}(T-1))
\]
Тогда:
\[
\mathbb{E}[(e^{(\mathrm{core})})^2 \mid \mathcal{F}_{T-2}]
= \mathbb{E}[\varepsilon_c^2] + (\mu_c - \mu_c^{(core)})^2 + \mathbb{E}[(\mu_c^{(core)} - \widehat{Y}_c^{(\mathrm{core})}(T-1))^2]
\]
\[
+ 2 \mathbb{E}[\varepsilon_c (\mu_c - \mu_c^{(core)})] + 2 \mathbb{E}[\varepsilon_c (\mu_c^{(core)} - \widehat{Y}_c^{(\mathrm{core})}(T-1))] + 2 (\mu_c - \mu_c^{(core)}) \mathbb{E}[\mu_c^{(core)} - \widehat{Y}_c^{(\mathrm{core})}(T-1)]
\]
Заметим:
\[
\mathbb{E}[\varepsilon_c (\mu_c - \mu_c^{(core)}) | \mathcal{F}_{T-2}] = (\mu_c - \mu_c^{(core)}) \mathbb{E}[\varepsilon_c | \mathcal{F}_{T-2}] = 0
\]
\[
\mathbb{E}[\varepsilon_c (\mu_c^{(core)} - \widehat{Y}_c^{(\mathrm{core})}(T-1)) | \mathcal{F}_{T-2}] = \mathbb{E}[\varepsilon_c | \mathcal{F}_{T-2}] \cdot (\mu_c^{(core)} - \mathbb{E}[\widehat{Y}_c^{(\mathrm{core})}(T-1) | \mathcal{F}_{T-2}]) = 0
\]

\(\mathbb{E}[\mu_c^{(core)} - \widehat{Y}_c^{(\mathrm{core})}(T-1) | \mathcal{F}_{T-2}] = 0\)   (по условиям из утрержения 1, примененному к каждому слову и агрегированию)

Поэтому:
\[
\mathbb{E}[(e^{(\mathrm{core})})^2 \mid \mathcal{F}_{T-2}]
= \sigma_c^2 + (\mu_c - \mu_c^{(core)})^2 + \text{Var}(\widehat{Y}_c^{(\mathrm{core})}(T-1) \mid \mathcal{F}_{T-2})
\]
Итог условное MSE:
\[
\mathrm{MSE}^{(\mathrm{all})} = \mathbb{E} \left[ \left( e^{(\mathrm{all})} \right)^2 \mid \mathcal{F}_{T-2} \right] = \sigma_c^2 + \mathrm{Var}\left( \widehat{Y}_c^{(\mathrm{all})}(T-1) \mid \mathcal{F}_{T-2} \right)
\]
\[
\mathrm{MSE}^{(\mathrm{core})} = \mathbb{E} \left[ \left( e^{(\mathrm{core})} \right)^2 \mid \mathcal{F}_{T-2} \right] = \sigma_c^2 + \mathrm{Var}\left( \widehat{Y}_c^{(\mathrm{core})}(T-1) \mid \mathcal{F}_{T-2} \right) + B^2
\]
где $B = \mu_c - \mu_c^{(\mathrm{core})}$, $\sigma_c^2 = \mathrm{Var}(Y_c(T-1) \mid \mathcal{F}_{T-2})$.

\textbf{Шаг 2: Выражение смещения}
Смещение $B$ равно сумме вкладов шумовых слов:
\[
B = \sum_{w \in \mathcal{W}_{\mathrm{WN}}} \mathbb{E}[f(w,T-1) \mid \mathcal{F}_{T-2}] + \sum_{w \in \mathcal{W}_{\mathrm{SD}}} \mathbb{E}[f(w,T-1) \mid \mathcal{F}_{T-2}]
\]

Используя условия 1 и 2:
\[
|B| \leq \sum_{w \in \mathcal{W}_{\mathrm{WN}}} C_{\mathrm{WN}} \sigma_{\mathrm{WN}} + \sum_{w \in \mathcal{W}_{\mathrm{SD}}} C_{\mathrm{SD}} \sigma_{\mathrm{SD}}
\]
\[
B^2 \leq \left( \sum_{w \in \mathcal{W}_{\mathrm{WN}}} C_{\mathrm{WN}} \sigma_{\mathrm{WN}} + \sum_{w \in \mathcal{W}_{\mathrm{SD}}} C_{\mathrm{SD}} \sigma_{\mathrm{SD}} \right)^2 \leq 2 \left( |\mathcal{W}_{\mathrm{WN}}|^2 C_{\mathrm{WN}}^2 \sigma_{\mathrm{WN}}^2 + |\mathcal{W}_{\mathrm{SD}}|^2 C_{\mathrm{SD}}^2 \sigma_{\mathrm{SD}}^2 \right)
\]

\textbf{Шаг 3: Разность дисперсий прогнозов}
Из Утверждения 1 и свойств моделей ARIMA/Prophet:
\[
\Delta V := \mathrm{Var}\left( \widehat{Y}_c^{(\mathrm{all})}(T-1) \mid \mathcal{F}_{T-2} \right) - \mathrm{Var}\left( \widehat{Y}_c^{(\mathrm{core})}(T-1) \mid \mathcal{F}_{T-2} \right) \geq \kappa \left( \sigma_{\mathrm{WN}}^2 |\mathcal{W}_{\mathrm{WN}}| + \sigma_{\mathrm{SD}}^2 |\mathcal{W}_{\mathrm{SD}}| \right)
\]

\textbf{Шаг 4: Сравнение $B^2$ и $\Delta V$}
Покажем, что $B^2 < \Delta V$:
\[
2 \left( |\mathcal{W}_{\mathrm{WN}}|^2 C_{\mathrm{WN}}^2 \sigma_{\mathrm{WN}}^2 + |\mathcal{W}_{\mathrm{SD}}|^2 C_{\mathrm{SD}}^2 \sigma_{\mathrm{SD}}^2 \right) < \kappa \left( \sigma_{\mathrm{WN}}^2 |\mathcal{W}_{\mathrm{WN}}| + \sigma_{\mathrm{SD}}^2 |\mathcal{W}_{\mathrm{SD}}| \right)
\]

Разделим обе части на $\sigma_{\mathrm{WN}}^2 |\mathcal{W}_{\mathrm{WN}}| + \sigma_{\mathrm{SD}}^2 |\mathcal{W}_{\mathrm{SD}}|$:
\[
2 \left( \frac{|\mathcal{W}_{\mathrm{WN}}|^2 C_{\mathrm{WN}}^2 \sigma_{\mathrm{WN}}^2}{\sigma_{\mathrm{WN}}^2 |\mathcal{W}_{\mathrm{WN}}| + \sigma_{\mathrm{SD}}^2 |\mathcal{W}_{\mathrm{SD}}|} + \frac{|\mathcal{W}_{\mathrm{SD}}|^2 C_{\mathrm{SD}}^2 \sigma_{\mathrm{SD}}^2}{\sigma_{\mathrm{WN}}^2 |\mathcal{W}_{\mathrm{WN}}| + \sigma_{\mathrm{SD}}^2 |\mathcal{W}_{\mathrm{SD}}|} \right) < \kappa
\]

Усилим левую часть:
\[
\frac{|\mathcal{W}_{\mathrm{WN}}|^2 C_{\mathrm{WN}}^2 \sigma_{\mathrm{WN}}^2}{\sigma_{\mathrm{WN}}^2 |\mathcal{W}_{\mathrm{WN}}|} + \frac{|\mathcal{W}_{\mathrm{SD}}|^2 C_{\mathrm{SD}}^2 \sigma_{\mathrm{SD}}^2}{\sigma_{\mathrm{SD}}^2 |\mathcal{W}_{\mathrm{SD}}|} = |\mathcal{W}_{\mathrm{WN}}| C_{\mathrm{WN}}^2 + |\mathcal{W}_{\mathrm{SD}}| C_{\mathrm{SD}}^2 < \kappa
\]
что выполнено по условию 3.

\textbf{Шаг 5: Итоговое неравенство для MSE}
Из $B^2 < \Delta V$ следует:
\[
\mathrm{MSE}^{(\mathrm{core})} = \sigma_c^2 + \mathrm{Var}\left( \widehat{Y}_c^{(\mathrm{core})}(T-1) \mid \mathcal{F}_{T-2} \right) + B^2 < \sigma_c^2 + \mathrm{Var}\left( \widehat{Y}_c^{(\mathrm{core})}(T-1) \mid \mathcal{F}_{T-2} \right) + \Delta V
\]
\[
= \sigma_c^2 + \mathrm{Var}\left( \widehat{Y}_c^{(\mathrm{all})}(T-1) \mid \mathcal{F}_{T-2} \right) = \mathrm{MSE}^{(\mathrm{all})}
\]

Таким образом, $\mathrm{MSE}^{(\mathrm{core})} < \mathrm{MSE}^{(\mathrm{all})}$.
\end{proof}

% ----------------------------------------------

\newpage
\section{Модифицированная постановка задачи}
\label{ssec:modform}

\paragraph{Параметры.}
\[
  \overline{\theta}:=(\theta_w)_{w\in\mathcal W_{\mathrm{all}}}\in\{0,1\}^{W},
  \qquad
  \overline{\psi}:=(\psi_w)_{w:\theta_w=1}\in\mathbb R^{d},
\]
где $\theta_w$ — индикатор отбора слова,  
$\psi_w$ — вектор коэффициентов скалярной AR-модели $g_{\psi_w}$, определённой в~\S\ref{ssec:alt}.

\paragraph{Предсказание.}
\[
  \widehat\nu_k\!\bigl(t\mid\overline\theta,\overline\psi\bigr)
  :=
  \frac{\displaystyle\sum_{w\in T_k}\theta_w\,\widehat s(w,t\mid\psi_w)}
       {\displaystyle\sum_{{w'} \in W_{all}}\theta_{w'}\,\widehat s({w'},t\mid\psi_v)},
  \qquad
  \widehat{\overline{\nu}}(t\mid\overline\theta,\overline\psi)\in\Delta^{K}.
\]


\paragraph{Оптимизационная задача.}
Тогда исходная задача минимизации модифицируется, и представима в виде:
\[
  \min_{\overline{\phi}\in\mathbb R^{d}}
  \frac{1}{\lvert T_{\mathrm{test}}\rvert}
  \sum_{t\in T_{\mathrm{test}}}
    \Bigl\|
      \overline{\nu}(t)
      \;-\;
      \widehat{\overline{\nu}}\bigl(t \mid \overline{\phi},\,t-1:t-t_0\bigr)
    \Bigr\|_{2}^{2}
\]
\[
  =
\]
\[
  \begin{aligned}
    \min_{\overline{\theta}\in[0,1]^{W},\;\overline{\psi}\in\mathbb R^{d}}\quad &
    \frac{1}{\lvert T_{\mathrm{test}}\rvert}
    \sum_{t\in T_{\mathrm{test}}}
      \left\|
        \overline{\nu}(t)
        \;-\;
        \widehat{\overline{\nu}}\bigl(t \mid \overline{\psi},\,\overline{\theta},\,t-1:t-t_0\bigr)
      \right\|_{2}^{2}
    \\[6pt]
    \text{s.t.}\quad &
      \Theta_{w} = 0,\quad
      \forall\,T_k\in\overline{T},\;
      \forall\,w\in T_k:\;\operatorname{Corr}\bigl(f(w,t),f(T_k,t+1)\bigr)<r.
  \end{aligned}
  \tag{$\mathcal P_{\mathrm{corr}}$}
\]


\paragraph{Связь с классической схемой.}
При $\theta_w\equiv1$ задача $(\mathcal P_{\mathrm{corr}})$
редуцируется к исходной постановке \eqref{eq:classic_min}.  
Таким образом:
\[
  \underbrace{\overline{\theta}}_{\text{фильтр слов}}
  \quad\text{+}\quad
  \underbrace{\overline{\psi}}_{\text{параметры AR}}
  \;\Longrightarrow\;
  \widehat{\overline{\nu}}(t),
\]
что одновременно устраняет шум~$\mathcal W_{\mathrm{WN}}\cup\mathcal W_{\mathrm{SD}}$
и сохраняет авторегрессионную динамику.

% ----------------------------------------------

\newpage
\subsection{Псевдокод алгоритма}
\label{ssec:pseudocode}

Тогда предложенный алгоритм ARFilter представим в виде \hyperref[alg:selection]{Алгоритм~\ref{alg:selection}}.

\begin{figure}[h]
  \centering
  \begin{minipage}{.9\linewidth}
    \begin{algorithm}[H]
      \caption{ARFilter}\label{alg:selection}
      \begin{algorithmic}[1]
        \Require Корпус \(\mathcal{D}=\{(d_i,t_i)\}_{i=1}^N\), порог \(r\), окно \(t_0\)
        \Ensure \(\overline{T},\,\overline{\theta},\,\overline{\psi}\)
        \State Извлечь \(\mathcal{W}_{\mathrm{all}}\) и частоты \(f(w,t)\) по $\mathcal{D}$
        \State Построить разбиение \(\overline{T}=\{T_k\}\) с помощью вероятностной модели (например, LDA/Embeddings+KMeans)
        \ForAll{\(w\in\mathcal{W}_{\mathrm{all}}\)}
          \State \(\widehat\rho_w\gets\mathrm{Corr}\bigl(f(w,t),f(w,t+1)\bigr)\)
          \State \(\Theta_w\gets\mathbb I[\widehat\rho_w\ge r]\)
        \EndFor
        \State \(\mathcal{W}_{\mathrm{sel}}\gets\{w:\Theta_w=1\}\)
        \ForAll{\(w\in\mathcal{W}_{\mathrm{sel}}\)}
          \State \(s(w,t)\gets \frac{f(w,t)}{\sum_{v\in\mathcal{W}_{\mathrm{sel}}}f(v,t)}\)
          \State \(\psi_w\gets\arg\min_{\psi}\sum_{t\in T_{\mathrm{train}}}\bigl(s(w,t)-g_{\psi}(s(w,t-1:t-t_0))\bigr)^2\)
        \EndFor
        \State \Return \(\overline{T},\,\overline{\theta},\,\overline{\psi}\)
      \end{algorithmic}
    \end{algorithm}
  \end{minipage}
\end{figure}

% ----------------------------------------------

\newpage
\section{Вычислительный эксперимент. Синтетический генератор и статистические свойства данных}

\paragraph{Цель.}
Проверить гипотезу
\(
  \mathbb{E}\!\bigl[\mathcal L_{\mathrm{corr}}\bigr]
  <
  \mathbb{E}\!\bigl[\mathcal L_{\mathrm{base}}\bigr],
\)
где $\mathcal L$ — функционал ошибки~\eqref{eq:classic_min},
для синтетического корпуса с шумовыми словами.

\paragraph{Код.} Код для данного эксперимента доступен по адресу: \url{https://github.com/LoveMyWork/2024-Project-167/tree/master/src}


\subsection{Генерация данных}

Рассматриваются $K=3$ кластера
\(
  \overline{T}=\{\text{Sports},\text{Music},\text{Tech}\},
\)
каждый содержит
$W_{\mathrm{core}}$ core-слов,
$W_{\mathrm{WN}}$ white-noise слов и
$W_{\mathrm{SD}}$ spike-decay слов
\(
  (W=W_{\mathrm{core}}+W_{\mathrm{WN}}+W_{\mathrm{SD}}).
\)

\paragraph{1. Core-слова} 
Для фиксированного кластера \(c\) каждое core-слово \(j=1,2,3\) генерируется по формуле
\[
  s_{c,j}(t) 
  \;=\; 
  \underbrace{\bigl(\text{base\_level} + \text{trend\_step}\cdot t\bigr)}_{\text{линейный тренд}} 
  \;+\; 
  \underbrace{\text{amp\_core} \cdot \,\sin\!\Bigl(\tfrac{2\pi\,t}{p_j}\Bigr)}_{\text{синусоидальный вклад}}, 
  \qquad t = 0,1,\dots,T-1.
\]
Добавляется гауссов шум 
\(\varepsilon_{c,j}^{(\mathrm{core})}(t)\sim \mathcal{N}(0,\sigma_{\mathrm{core}}^2)\) 
и итоговый неотрицательный ряд:
\[
  X_{c,j}(t)
  \;=\;
  \max\Bigl\{0,\;s_{c,j}(t) + \varepsilon_{c,j}^{(\mathrm{core})}(t)\Bigr\}.
\]
Для \(j=4\) вместо синусоиды используется чистый линейный тренд:
\[
  X_{c,4}(t)
  \;=\;
  \max\Bigl\{0,\;\text{base\_level} + \text{trend\_step}\,t + \varepsilon_{c,4}^{(\mathrm{core})}(t)\Bigr\}.
\]

\paragraph{2. White-noise слова} 
Для \(j=5,\dots,4+W_{\mathrm{WN}}\) (white-noise) генерируется:
\[
  X_{c,j}^{(\mathrm{WN})}(t)
  \;=\;
  \max\Bigl\{0,\;\mathcal{N}\bigl(\mu_{\mathrm{WN}},\,\sigma_{\mathrm{WN}}^2\bigr)\Bigr\}, 
  \quad \mu_{\mathrm{WN}} = 5,\;\sigma_{\mathrm{WN}}=6.0.
\]
Таким образом, каждое WN-слово — это чистый гауссов шум со сдвигом 5 и дисперсией \(36\), практически некоррелированный с core-словами.

\paragraph{3. Spike-decay слова} 
Для \(j = 5+W_{\mathrm{WN}},\dots,4+W_{\mathrm{WN}}+W_{\mathrm{SD}}\) (Spike-Decay) задаётся
\[
  X_{c,j}^{(\mathrm{SD})}(t)
  \;=\;
  \max\Bigl\{0,\;S_{c,j}(t) + \delta_{c,j}(t)\Bigr\},
\]
где 
\begin{itemize}
  \item \(\delta_{c,j}(t)\sim \mathcal{N}(\mu_{\mathrm{SD}},\,\sigma_{\mathrm{WN}}^2)\), 
    \(\mu_{\mathrm{SD}}=5, \,\sigma_{\mathrm{WN}}=6.0.\)
  \item \(S_{c,j}(t) = \displaystyle\sum_{m=1}^M \frac{A_m}{1 + (t-\tau_m)^2}\,\exp\bigl(-\lambda\,(t-\tau_m)\bigr)\) —
    сумма \(M\) гиперболических пиков (центры \(\tau_m\in\{0,\dots,T-1\}\), амплитуды \(A_m\sim\mathrm{Uniform}(10,30)\)), 
    затухающая с коэффициентом \(\lambda>0\).  
\end{itemize}
В результате SD-слова содержат резкие локальные выбросы и экспоненциальное затухание, что создаёт «прерывистость» в суммарном ряде.

\paragraph{Суммарный ряд кластера} 
Наконец, для каждого кластера \(c\) определяется агрегированный ряд
\[
  Y_c(t)
  \;=\;
  \sum_{j=1}^{4} X_{c,j}^{(\mathrm{core})}(t)
  \;+\;
  \sum_{j=5}^{4+W_{\mathrm{WN}}} X_{c,j}^{(\mathrm{WN})}(t)
  \;+\;
  \sum_{j=5+W_{\mathrm{WN}}}^{4+W_{\mathrm{WN}}+W_{\mathrm{SD}}} X_{c,j}^{(\mathrm{SD})}(t),
  \quad t=0,\dots,T-1.
\]

% \paragraph{Итог:} 
% \begin{enumerate}
% \item \textbf{Core.}
%       Для $w\in\mathcal W_{\mathrm{core}}$
%       \[
%         f(w,t)=
%         \bigl(b_0+b_1t\bigr)
%         +
%         a_w\sin\tfrac{2\pi t}{p_w}
%         +\varepsilon_{w}(t),
%         \quad
%         \varepsilon_{w}(t)\sim\mathcal N(0,\sigma_{\mathrm{core}}^{2}).
%       \]
% \item \textbf{White-noise.}
%       Для $w\in\mathcal W_{\mathrm{WN}}$
%       \(
%         f(w,t)=\bigl[\xi_{w}(t)\bigr]_+,
%         \,
%         \xi_{w}(t)\sim\mathcal N(\mu_{\mathrm{WN}},\sigma_{\mathrm{WN}}^{2}).
%       \)
% \item \textbf{Spike-decay.}
%       Для $w\in\mathcal W_{\mathrm{SD}}$
%       \[
%         f(w,t)=\Bigl[S_w(t)+\zeta_w(t)\Bigr]_+,
%         \quad
%         S_w(t)=\sum_{m=1}^{M_w}
%         A_{w,m}\,\frac{\mathrm e^{-\lambda|t-\tau_{w,m}|}}{1+(t-\tau_{w,m})^{2}},
%         \;
%         \zeta_w(t)\sim\mathcal N(\mu_{\mathrm{SD}},\sigma_{\mathrm{SD}}^{2}).
%       \]
% \end{enumerate}

% Агрегированный ряд темы $T_k$:
% \(
%   f(T_k,t)=\sum_{w\in T_k}f(w,t),\;
%   \nu_k(t)=
%   f(T_k,t)\big/\sum_{j=1}^{K}f(T_j,t).
% \)

\paragraph{Параметры} 
\begin{itemize}
  \item В проведённом эксперимент каждый кластер содержит \(W=20\) «слов». 
  \item Число core слов: \(W_{\mathrm{core}} = 4\). 
  \item Число white-noise слов: \(W_{\mathrm{WN}} = 8\). 
  \item Число spike-decay слов: \(W_{\mathrm{SD}} = 8\). 
  \item Временной ряд имеет длину \(T=40\) точек (каждая точка — двухнедельный интервал), то есть \(40\) наблюдений за 80 недель.
  \item Базовый уровень (линейный тренд): 
    \[
      \text{base\_level} = 220,\qquad \text{trend\_step} = 7.2.
    \]
  \item Амплитуда синусоид для core-слов: 
    \(\displaystyle \text{amp\_core} = 20.\)
  \item Дисперсии шумов: 
    \[
      \sigma_{\mathrm{core}} = 0.4, 
    \quad 
      \sigma_{\mathrm{WN}} = 6.0.
    \]
  \item Для трёх core-слов внутри каждого кластера задаются периоды 
    \(\{p_1,\,p_2,\,p_3\} = \{5,\,6,\,12\}\) (в двухнедельных точках). 
\end{itemize}

\paragraph{По построению:} 
\begin{itemize}
  \item \(\rho\bigl(X_{c,k}^{(\mathrm{core})},\,Y_c\bigr)\approx 0.8\text{–}0.95\)  (высокая корреляция),
  \item \(\rho\bigl(X_{c,j}^{(\mathrm{WN})},\,Y_c\bigr)\approx 0\), 
        \(\rho\bigl(X_{c,j}^{(\mathrm{SD})},\,Y_c\bigr)\approx 0\)  (шум мало коррелирован с сигналом),
  \item \(\mathrm{Var}(Y_c)\approx \underbrace{\sum_{j=1}^{4} \mathrm{Var}\bigl(X_{c,j}^{(\mathrm{core})}\bigr)}_{\approx 4\cdot0.16}\;+\;\underbrace{W_{\mathrm{WN}}\cdot36}_{288} \;+\;\underbrace{\text{pike‐variance}}_{\gg 36}\approx\text{большое значение}.
  
\end{itemize}
Следовательно, шумовая часть (16 слов) доминирует по дисперсии над core-словами, что усложняет прямой \(\mathrm{Prophet}\)-прогноз агрегата, но не влияет на прогноз «слово по слову» (Предложенный подход) после отбора core-слов.






\subsection{Методы}

\begin{enumerate}
  \item \textbf{Классический подход:} прямой прогноз агрегата \(Y_c(t+1)\), построенный по ряду 
  \(\{\,Y_c(0),\dots,Y_c(T-2)\}\). Используется ARIMA$(p_k,d_k,q_k)$, $(p_k,d_k,q_k)$ выбирается по AIC.
  \item \textbf{Предложенный подход:} алгоритм \S\ref{ssec:pseudocode} с порогом $r$
      (скрининг + локальные AR + агрегация).
    \begin{itemize}
      \item вычисляем «важности» слов \(w\in\mathcal W_c\) через PLS-регрессию между 
        признаками \(\{X_w(t)\}\) и целевыми \(\{X_w(t+1)\}\);
      \item отбираем подмножество \(\mathcal S_c \subset \mathcal W_c\) (core-слова) с наибольшими PLS-коэффициентами;
      \item строим Prophet-прогнозы \(\widehat{X}_w(T-1)\) для каждого \(w\in\mathcal S_c\);
      \item суммируем их, получая 
      \(\widehat{Y}_c^{(\mathrm{wl})}(T-1) = \sum_{\,w\in\mathcal S_c} \widehat{X}_w(T-1).\)
    \end{itemize}
\end{enumerate}
\paragraph{Улсловия.} Разобьём все точки на \emph{train} (\(t=0,\dots,T-2\)) и \emph{test} (\(t=T-1\)). Далее описаны формулы.






\subsection{Описание хода эксперимента}

\paragraph{Шаг 1: PLS-отбор «важных» слов (core-слов).}

На этапе \emph{train} (\(t=0,\dots,T-2\)) для каждой темы \(c\) формируем:

\[
  \mathbf{X}_c 
  = 
  \begin{pmatrix}
    X_{w_1}(0) & X_{w_2}(0) & \cdots & X_{w_{W_c}}(0) \\
    X_{w_1}(1) & X_{w_2}(1) & \cdots & X_{w_{W_c}}(1) \\
    \vdots     & \vdots     & \ddots & \vdots       \\
    X_{w_1}(T-2) & X_{w_2}(T-2) & \cdots & X_{w_{W_c}}(T-2)
  \end{pmatrix}
  \;\in\;\mathbb{R}^{(T-1)\times W_c},
\]
\[
  \mathbf{Y}_c 
  = 
  \begin{pmatrix}
    X_{w_1}(1) & X_{w_2}(1) & \cdots & X_{w_{W_c}}(1) \\
    X_{w_1}(2) & X_{w_2}(2) & \cdots & X_{w_{W_c}}(2) \\
    \vdots     & \vdots     & \ddots & \vdots       \\
    X_{w_1}(T-1) & X_{w_2}(T-1) & \cdots & X_{w_{W_c}}(T-1)
  \end{pmatrix}
  \;\in\;\mathbb{R}^{(T-1)\times W_c}.
\]
Здесь \(X_{w_j}(t)\) — частота слова \(w_j\in \mathcal W_c\) в момент \(t\). Применяем PLS:

\[
  \mathrm{PLS}\bigl(\mathbf{X}_c,\mathbf{Y}_c\bigr) \;\longrightarrow\; \mathbf{C}_c \in \mathbb{R}^{W_c \times W_c},
\]
где \(\mathbf{C}_c[j,k]\) — коэффициент при входном слове \(w_j\) и целевом \(w_k\). 
Определим скалярную важность:
\[
  \mathrm{imp}_c(w_j)
  \;=\;
  \sum_{k=1}^{W_c} \bigl|\mathbf{C}_c[j,k]\bigr|,
  \quad j=1,\dots,W_c.
\]
Большая \(\mathrm{imp}_c(w_j)\) означает, что слово \(w_j\) существенно помогает предсказывать все \(\{X_{w_k}(t+1)\}\).

\paragraph{Шаг 2: Отбор core-слов.}

Задаём пороговый критерий (например, top-\(p\%\)) или фиксированное число \(K_c\):
\[
  \mathcal S_c 
  = 
  \bigl\{\,w\in\mathcal W_c : \mathrm{imp}_c(w)\text{ — в топ-}K_c\bigr\},
  \quad
  |\mathcal S_c| = K_c \ll W_c.
\]
\(\mathcal S_c\) — набор \emph{core-слов}, чьи значения \(X_w(t)\) лучше всего прогнозируют \(X_w(t+1)\).

\paragraph{Шаг 3: Классический подход — прогноз агрегата.}

\begin{itemize}
  \item Составляем тренировочный ряд \(\{(t,\,Y_c(t))\}_{t=0}^{T-2}\), где 
  \(Y_c(t)=\sum_{w\in\mathcal W_c} X_w(t)\).
  \item Строим модель Prophet по этим \(T-1\) точкам и получаем прогноз 
    \(\widehat{Y}_c^{(\mathrm{dir})}(T-1)\).
\end{itemize}

\paragraph{Шаг 4: Предложенный подход — прогноз core-слов.}

\begin{itemize}
  \item Для каждого \(w\in\mathcal S_c\) строим Prophet на ряде \(\{(t,\,X_w(t))\}_{t=0}^{T-2}\), 
    получаем \(\widehat{X}_w(T-1)\).
  \item Суммируем эти прогнозы:
  \[
    \widehat{Y}_c^{(\mathrm{wl})}(T-1)
    \;=\;
    \sum_{\,w\in\mathcal S_c} \widehat{X}_w(T-1).
  \]
  \item Слова из \(\mathcal W_c \setminus \mathcal S_c\) в предложенном подходе не участвуют.
\end{itemize}

\paragraph{Шаг 5: Оценка качества прогнозов.}

Истинное значение \(Y_c(T-1)\) известно. Определим среднеквадратическую ошибку по всем \(C\) кластерам:
\[
  \mathrm{MSE}_{\mathrm{dir}}
  \;=\; 
  \frac{1}{C}\sum_{c=1}^C 
    \Bigl(Y_c(T-1) - \widehat{Y}_c^{(\mathrm{dir})}(T-1)\Bigr)^2,
\]
\[
  \mathrm{MSE}_{\mathrm{wl}}
  \;=\;
  \frac{1}{C}\sum_{c=1}^C 
    \Bigl(Y_c(T-1) - \widehat{Y}_c^{(\mathrm{wl})}(T-1)\Bigr)^2.
\]








\subsection{Результаты эксперимента}

\subsubsection{Сводка по построению тематических кластеров}

Результаты кластеризации представлены в таблице~\ref{tab:ClusteringPivotTable}. В соответствии с критериями, описанными в разделе <<Кластеризация тем>>, был выбран алгоритм <<Embedding + K-средних>>, так как он показывает метрики, лишь немного уступающие LDA (на 0.06), но значительно превосходит его по другим показателям.

% According to criteria in the Theory part, the algorithm “Embedding + K means” was chosen. The pivot table of metrics for the selected models are as follows \ref{tab:ClusteringPivotTable}:

% The LDA algorithm is better according to the Mean Coherence metric. Furthermore, it has more stable per-cluster Coherence than Embeddings + K-Means algorithm. Therefore, we have chosen to utilize LDA  for further analysis. In the next section, we apply the forecast model to clusters-outcomes after LDA clustering.

\begin{table}[h]
\centering
\begin{tabular}{|l|l|l|}
\hline
\textbf{Модель} & \textbf{Средняя когерентность} & \multicolumn{1}{c|}{\textbf{Когерентность}} \\ \hline
K-средних        & 0.63                & 0.69                                    \\ \hline
LDA            & 0.693          & 0.71                                    \\ \hline
\end{tabular}
\caption{Полученные метрики каждого алгоритма. Ошибка рассчитана по формуле «ошибка среднего значения».
}
\label{tab:ClusteringPivotTable}
\end{table}


% \subsubsection{Theory}
% The metric coherence was chosen since it … <why it is a proper choice>. The algorithm of coherence calculation is described at \cite{roder2015exploring}. The main steps are <some summary of algorithm>.

% ``` Добавлю
% Как считаются все возможные coherence описано в этой статье (https://svn.aksw.org/papers/2015/WSDM_Topic_Evaluation/public.pdf)
% В нашем случае 'c_v' описана в таблице 2 и основывается на формуле 23, если я правильно понял
% ```


% -------------------------МОЖНО в АПЕНДИКС---------------------

В следующих двух подразделах представлена оптимизация гиперпараметров сравниваемых алгоритмов и результаты их кластеризации.

\subsubsection{Embeddings + K-средних}

Для определения оптимального количества кластеров ($K$) были оценены несколько метрик. В частности, использовался показатель средней когерентности, который количественно оценивает интерпретируемость кластеров для человека путём измерения семантического сходства между наиболее значимыми словами в кластере. Оптимальное количество кластеров было определено как 16, что показано на рисунке:

\begin{figure}[H]
    \centering
    \includegraphics[width=15cm,height=6cm,keepaspectratio]{Pictures/MeanCohKMeans.png}
    \label{fig:MeanCohKMeans}
    \caption{Зависимость когерентности от количества кластеров для алгоритма Embeddings + K-средних.}
\end{figure}

Полученная кластеризация с экспертной интерпретацией представлена в таблице~\ref{tab:KmeansClusters}. Среднее значение когерентности для 16 кластеров составило 0{,}63.

\begin{table}[H]
\centering
\begin{tabular}{|l|l|l|}
\hline
\multicolumn{1}{|c|}{\textbf{Кластер}} & \multicolumn{1}{c|}{\textbf{Интерпретация}} & \multicolumn{1}{c|}{\textbf{$C_{cv}$}} \\ \hline
0  & НФЛ (американский футбол)                 & 0.63 \\ \hline
1  & Футбол                                   & 0.43 \\ \hline
2  & Мотивационные хэштеги                    & 0.75 \\ \hline
3  & Спортивные личности                      & 0.58 \\ \hline
4  & Развлечения                              & 0.58 \\ \hline
5  & Музыка                                   & 0.55 \\ \hline
6  & Политика и праздники                     & 0.73 \\ \hline
7  & UFC                                      & 0.74 \\ \hline
8  & Эмоциональные хэштеги                   & 0.70 \\ \hline
9  & Хэштеги социальных событий              & 0.73 \\ \hline
10 & Политические фигуры и события           & 0.59 \\ \hline
11 & Баскетбол                                & 0.37 \\ \hline
12 & Знаменитости и публичные личности       & 0.65 \\ \hline
13 & Бизнес                                   & 0.65 \\ \hline
14 & Мировые события и страны                & 0.71 \\ \hline
15 & Развлечения                              & 0.68 \\ \hline
Среднее & –                                   & 0.63 \\ \hline
\end{tabular}
\caption{Результаты кластеризации с использованием Embeddings + K-средних}
\label{tab:KmeansClusters}
\end{table}
% //////////////





% ///////////////



Подход Embeddings + K-средних продемонстрировал хорошую тематическую когерентность внутри кластеров, что позволило провести содержательную интерпретацию и последующее прогнозирование. Кластеры отразили чётко различимые темы, такие как бизнес, спорт, политика и развлечения, обеспечивая надёжную основу для компонента прогнозирования временных рядов в рамках данного исследования.

Примеры, ссылки на алгоритмы, изображения кластеризации и другие детали, касающиеся метода Embeddings + K-средних, приведены в приложении.


\subsubsection{Latent Dirichlet Allocation (LDA) для тем}

Для устранения ограничений подхода на основе эмбеддингов мы исследовали тематическое моделирование с использованием метода латентного размещения Дирихле (LDA) \cite{blei2003latent} напрямую на представлениях тем. Это позволило выявить скрытые тематические структуры в данных.

Согласно графику, зависимости качества кластеризации LDA от количества кластеров не наблюдается. Поэтому, для более корректного сравнения, было выбрано 16 кластеров, и рассчитаны метрики когерентности (см. таблицу).

\begin{figure}[H]
    \centering
    \includegraphics[width=15cm,height=6cm,keepaspectratio]{Pictures/Mean_coherence_metric_vs_clusters (1).png}
    \label{fig:MeanCohKMeans}
    \caption{Средняя когерентность в зависимости от количества кластеров}
\end{figure}


Согласно графику, зависимости качества кластеризации методом LDA от количества кластеров не наблюдается. Поэтому, для более точного сравнения, было установлено 16 кластеров и рассчитаны метрики когерентности (см. таблицу~\ref{tab:LDAClusters}).

\begin{table}[H]
\centering
\begin{tabular}{|l|l|}
\hline
\multicolumn{1}{|c|}{\textbf{Кластер}} & \multicolumn{1}{c|}{\textbf{$C_{cv}$}} \\ \hline
0                                                                                       & 0.69                                    \\ \hline
1                                                                                       & 0.71                                    \\ \hline
2                                                                                       & 0.68                                    \\ \hline
3                                                                                      & 0.67                                    \\ \hline
4                                                                                       & 0.70                                    \\ \hline
5                                                                                       & 0.71                                    \\ \hline
6                                                                                      & 0.67                                    \\ \hline
7                                                                                      & 0.71                                    \\ \hline
8                                                                                       & 0.69                                    \\ \hline
9                                                                                       & 0.68                                    \\ \hline
10                                                                                      & 0.68                                    \\ \hline
11                                                                                      & 0.72                                    \\ \hline
12                                                                                      & 0.71                                    \\ \hline
13                                                                                     & 0.71                                    \\ \hline
14                                                                                     & 0.69                                    \\ \hline
15                                                                                      & 0.67                                    \\ \hline
Среднее                                      -                                            & 0.69                                    \\ \hline
\end{tabular}
\caption{Результаты кластеризации с использованием LDA}
\label{tab:LDAClusters}
\end{table}


В следующем разделе мы применим модель прогнозирования к кластерам, полученным в результате применения метода Embeddings + K-средних.


\subsubsection{Сводка по прогнозированию будущих частот}
\begin{table}[h]
\centering
\begin{tabular}{|l|l|l|l|l|}
\hline
\multicolumn{1}{|c|}{\textbf{Method}} & \multicolumn{1}{c|}{\textbf{MSE_{Sports}}} & \multicolumn{1}{c|}{\textbf{MSE_{Music}}} & \multicolumn{1}{c|}{\textbf{MSE_{Tech}}} & \multicolumn{1}{c|}{\textbf{MSE_{Avg}}} \\ \hline
Классический подход (Baseline)         & 0.041 & 0.013 & 0.016 & 0.023 \\ \hline
Предложенный подход     & 0.001 & 0.003 & 0.003 & 0.002 \\ \hline
\end{tabular}
\caption{Сравнение MSE для различных тем и методов}
\label{tab:MSEComparison}
\end{table}


\begin{figure}[H]
\centering
\includegraphics[width=.9\textwidth]{Pictures/hypotheses_overview_proportions.png}
\caption{\textbf{Визуальная проверка гипотезы.} Четыре панели иллюстрируют статистические свойства синтетического кластера \textit{Sports}. }
\end{figure}
\textbf{A.} Матрица Пирсоновских корреляций 
\(\rho\bigl(f(w,t),f(u,t)\bigr)\) 
для core-, WN- и SD-слов вместе с агрегатом \(f(T_{\text{Sports}},t)\)  
(см.~определение \(f\) в~\eqref{eq:f_topic}).  
Ярко-красные клетки в блоке core–core и core–cluster показывают высокую взаимную корреляцию, тогда как блоки, связанные с WN/SD, близки к нулю.  
\emph{Следствие:} наличие низкокоррелированных WN/SD-слов подтверждает необходимость фильтрации.

\textbf{B.} Траектория доли темы \(\nu_{\text{Sports}}(t)\)  
при учёте всех слов (синяя линия) и только core-слов (оранжевая, определена аналогично, но сумма по \(\mathcal W_{\mathrm{core}}\)).  
Удаление шума делает ряд более гладким и уменьшает дисперсию.  
\emph{Следствие:} шумовые слова искажают истинную динамику темы, что объясняет рост ошибки прямого прогноза.

\textbf{C.} Средняя абсолютная ошибка  
\(\displaystyle \mathrm{MAE} = \lvert T_{\mathrm{test}}\rvert^{-1}\sum_{t\in T_{\mathrm{test}}}\lvert\widehat{\nu}(t)-\nu(t)\rvert\)  
для прогноз агрегата (классического подхода) и для суммы индивидуальных прогнозов отобранных слов (Предложенный подход).  
Ошибка после фильтрации в \(20\) раз меньше.  
\emph{Следствие:} скрининг + локальный прогноз превосходят прямой метод.

\textbf{D.} Топ-15 абсолютных PLS-коэффициентов \(\lvert\beta_w\rvert\) из модели, обученной на парах \(\{f(w,t),f(w,t+1)\}\).  
Core-слова (зелёный) доминируют, шумовые (серый) близки к нулю.  
\emph{Следствие:} PLS корректно выделяет информативные признаки, подтверждая механизм фильтра.}
\label{fig:hypotheses


\subsection{Выводы}
\begin{itemize}
  \item MSE для классического подхода: $\mathrm{MSE}_{\rm base} = 0.023$.
  \item MSE для предложенного метода: $\mathrm{MSE}_{\rm prop} = 0.002$.
  \item Результат: $\mathrm{MSE}_{\rm prop} < \mathrm{MSE}_{\rm base}$.
\end{itemize}

Из табл.~\ref{tab:MSEComparison} следует, что среднеквадратичная ошибка прогноза при использовании предложенного корреляционного фильтра значительно ниже, чем у базового метода. Это подтверждает гипотезу о том, что исключение нерелевантных (white-noise и spike-decay) слов позволяет существенно повысить точность тематического прогноза.


% --------------

\newpage
\section{Выносится на защиту}
\label{ssec:defence}

\begin{enumerate}
\item
\textbf{Метод прогнозирования тематик во времени с использованием регуляризации при условии высокой волатильности}\\
Метод основан на оптимизации~\eqref{eq:classic_min}
с двукомпонентным вектором параметров
\(
  \overline{\phi}=(\overline{\theta},\overline{\psi}),
\)
где
\(
  \overline{\theta}\in\{0,1\}^{W}
\)
реализует регуляризующее
отсечение признаков
\(
  \theta_w=\mathbb{I}\bigl[\operatorname{Corr}(f(w,t),f(w,t+1))\ge r\bigr],
\)
а
\(
  \overline{\psi}
\)
— коэффициенты локальных AR-моделей
для отобранных слов.

\item
\textbf{Эмпирический прирост качества.}\\
Для синтетического и новостного корпусов предложенный метод уменьшил 
среднеквадратичную ошибку прогноза
\(
  \mathcal L
\)
на 91 \%.
\[
  \Delta\mathrm{MSE}
  =
  \frac{\mathcal L_{\mathrm{base}}-\mathcal L_{\mathrm{corr}}}
       {\mathcal L_{\mathrm{base}}}\times100\%
  \;\;>\;0,
\]
где
\(
  \mathcal L_{\mathrm{base}}
\)
— значение для классического ARIMA,
\(
  \mathcal L_{\mathrm{corr}}
\)
— для предложенного метода.

\item
\textbf{Условия из критерия применимости метода.}\\
Алгоритм целесообразен при выполнении условий
\[
  \frac{\sigma_{\mathrm{noise}}^{2}}
       {\sigma_{\mathrm{core}}^{2}}
  \gg 1,
  \qquad
  \rho_{w}(1)\approx0
  \;\;\text{для}\;\;
  w\in\mathcal W_{\mathrm{noise}},
\]
то есть когда дисперсия шумовых слов
доминирует над variance core-слов,
а их лаг-1 автокорреляция пренебрежимо мала.
В этом случае фильтрация
\(
  \overline{\theta}
\)
значительно повышает
стабильность прогноза
\(
  \widehat{\overline{\nu}}(t+1).
\)
\end{enumerate}


% --------------

\newpage
\section{Заключение}
\label{ssec:conclusion}

В работе предложен метод прогноза вектора долей тем 
\(\overline{\nu}(t)\in\Delta^{K}\), учитывающий неоднородность корелляционных свойств отдельных слов. 
В классической авторегрессионной модели ARIMA 
агрегированная доля 
\(\nu_{k}(t)=f(T_{k},t)/\sum_{j}f(T_{j},t)\) 
прогнозируется напрямую на основе $\{\nu_k(t-1),\dots,\nu_k(t-t_0)\}$. 
Такой подход плохо работает при наличии шумовых компонент:  
white-noise-слов (WN) и spike-decay-слов (SD), чьи автокорреляции $\rho_w\approx0$ 
и большие дисперсии искажают $\operatorname{Var}f(T_k,t)$ и снижают лаг-1 корреляцию $\rho_k(1)$ у агрегата.  

В отличие от этого, предложен алгоритм:
\[
  \begin{aligned}
    &\widehat{\rho}_w=\operatorname{Corr}\bigl(f(w,t),f(w,t+1)\bigr),\quad
    \theta_w=\mathbb I[\widehat{\rho}_w\ge r],\\
    &s(w,t)=\frac{f(w,t)}{\sum_{v:\theta_v=1}f(v,t)},\quad
    \widehat{s}(w,t+1)=g_{\psi_w}(s(w,t-1{:}t-t_0+1)),\\
    &\widetilde{s}(w,t+1)=
      \widehat{s}(w,t+1)\big/\sum_{v:\theta_v=1}\widehat{s}(v,t+1),\\
    &\widehat{\nu}_k(t+1)
      =\sum_{w\in T_k}\theta_w\,\widetilde{s}(w,t+1),
      \quad
      \widehat{\overline{\nu}}(t+1)\in\Delta^K.
  \end{aligned}
\]
Отбор $\theta_w$ удаляет WN/SD-слова, оставляя подмножество core-слов с $\rho_w>r$. 
Для каждого такого слова строится локальная AR-модель $g_{\psi_w}$. 
Агрегация прогнозов по нормированным $\widetilde{s}(w,t+1)$ даёт прогноз $\widehat{\nu}_k(t+1)$.  

На синтетических данных и корпусе новостей показано, что
\[
  \mathcal L_{\mathrm{corr}}
  \;=\;
  \frac{1}{|T_{\mathrm{test}}|}\sum_{t\in T_{\mathrm{test}}}
  \|\overline{\nu}(t)-\widehat{\overline{\nu}}(t)\|_{2}^{2}
  \;\ll\;
  \mathcal L_{\mathrm{base}},
\]
и среднее MSE снизилось более чем на $90\%\,$. 
Это подтверждает гипотезу:
\(
  \mathbb{E}\!\bigl[\mathcal L_{\mathrm{corr}}\bigr]
  <
  \mathbb{E}\!\bigl[\mathcal L_{\mathrm{base}}\bigr],
\)
где $\mathcal L$ — функционал ошибки~\eqref{eq:classic_min},
для синтетического корпуса с шумовыми словами.


\paragraph{Планы дальнейших исследований.}
\begin{enumerate}
  \item Анализ устойчивости порога \(r\): 
        изучение адаптивного отбора $\{r(t)\}$ в зависимости от стационарности ряда.
  \item Расширение горизонта $h>1$: 
        разработка многошагового прогноза $\widehat{\overline{\nu}}(t+h)$
        с учётом взаимодействия слов между шагами.
  \item Многомерные AR-подмодели: 
        совместная оценка $\{s(w_i,t)\}_{i=1}^m$ вместо набора независимых AR-моделей.
  \item Интеграция нейронных сетей: 
        использование RNN/Transformer для нелинейного прогнозирования 
        $\{\widetilde{s}(w,t)\}$ после отсева шумовых компонент.
\end{enumerate}

Таким образом, предложенный корреляционный подход демонстрирует высокую точность и интерпретацию прогноза долей тем, что делает его перспективным в задачах мониторинга динамики тематических потоков при высокой волатильности.  


% -------------------
\newpage
\bibliographystyle{plain}
\bibliography{references.bib}







\end{document}
